\documentclass{article}
\usepackage[spanish,es-nodecimaldot,es-noquoting]{babel}
\usepackage[letterpaper,includehead,top=2cm,bottom=3cm,left=3cm,right=3cm,headheight=70pt,headsep=13pt,footskip=20pt]{geometry}
\usepackage{amsmath,amssymb,graphicx,fancyhdr,multicol}
\usepackage{booktabs,longtable,tabularx,array,xcolor,listings,tikz}
\usetikzlibrary{arrows.meta,positioning,calc}
\usepackage{xurl}
\usepackage[colorlinks=true,allcolors=blue!55!black,pdfencoding=auto]{hyperref}
\usepackage{caption}
\captionsetup{font=small,labelfont=bf,hypcap=false}
\hypersetup{pdftitle={LogiMarket 360 — Informe APE 04},pdfsubject={Arquitectura, patrones GoF y evidencias de aceptación}}
\definecolor{acento}{HTML}{12475B}
\definecolor{fondo}{HTML}{F2F6F8}

\definecolor{lgmuted}{HTML}{374E58}
\definecolor{lgline}{HTML}{708998}
\definecolor{lgerror}{HTML}{8D2926}
\tikzset{
  lgBox/.style={draw=lgline,fill=fondo,rounded corners=1.5pt,align=center,
    minimum height=8.5mm,inner sep=3.0pt,font=\small,text=black!91},
  lgHeading/.style={font=\small\bfseries,text=acento,align=center},
  lgUse/.style={lgBox,rounded corners=11pt,text width=4.6cm},
  lgNote/.style={font=\footnotesize,text=lgmuted,align=center,text width=14.2cm},
  lgLabel/.style={font=\footnotesize,text=black!87,fill=white,align=center,inner sep=1.2pt},
  lgParticipant/.style={lgBox,font=\footnotesize\bfseries,draw=acento,minimum height=10mm},
  lgLife/.style={draw=lgline,densely dotted,line width=.56pt},
  lgActive/.style={draw=acento,fill=teal!14,line width=.55pt},
  lgBand/.style={font=\footnotesize\bfseries,draw=lgline,fill=fondo,rounded corners=1pt,
    inner sep=3pt,text width=12.7cm,align=left,text=acento},
  lgCall/.style={-{Latex[length=2.1mm]},draw=acento,line width=.64pt,rounded corners=2pt},
  lgReturn/.style={-{Latex[length=1.8mm]},draw=lgmuted,dashed,line width=.63pt},
  lgError/.style={-{Latex[length=2mm]},draw=lgerror,dash pattern=on 3pt off 2pt,line width=.85pt,rounded corners=2pt},
  lgAssoc/.style={draw=acento,line width=.64pt,rounded corners=2pt},
  lgDep/.style={-{Stealth[open,length=2mm]},draw=acento,dashed,line width=.64pt,rounded corners=2pt},
  lgInherit/.style={-{Triangle[open,length=2.6mm,width=2.25mm,fill=white]},draw=acento,line width=.65pt,rounded corners=2pt},
  lgReal/.style={lgInherit,dashed},
  lgComp/.style={{Diamond[length=2.6mm,width=1.9mm]}-,draw=acento,line width=.65pt,rounded corners=2pt}
}
\newsavebox{\lgDiagram}
\lstset{language={[Sharp]C},basicstyle=\ttfamily\footnotesize,
  frame=single,rulecolor=\color{black!22},backgroundcolor=\color{fondo},
  numbers=left,numberstyle=\scriptsize\color{black!65},numbersep=6pt,
  breaklines=true,breakatwhitespace=true,showstringspaces=false,columns=fullflexible,
  xleftmargin=1.8em,framexleftmargin=1.25em,keepspaces=true,
  literate={á}{{\'a}}1 {é}{{\'e}}1 {í}{{\'i}}1 {ó}{{\'o}}1 {ú}{{\'u}}1 {ñ}{{\~n}}1 {Á}{{\'A}}1 {É}{{\'E}}1 {Ó}{{\'O}}1}
\newcommand{\plabel}[1]{\par\smallskip\noindent\textbf{#1}\quad\ignorespaces}
\newcommand{\vistafig}[4]{\par\smallskip\noindent\begin{minipage}{\linewidth}\centering
  \includegraphics[width=\linewidth,height=#3,keepaspectratio]{capturas/#1}
  \captionof{figure}{#4}\label{fig:#2}\end{minipage}\par\medskip}
\newcommand{\wirefig}[4]{\par\smallskip\noindent\begin{minipage}{\linewidth}\centering\includegraphics[width=\linewidth,height=#3,keepaspectratio]{wireframes/#1}\captionof{figure}{#4}\label{fig:#2}\end{minipage}\par\medskip}
\newcommand{\diagrama}[3]{\par\medskip\noindent\begin{minipage}{\linewidth}\centering
  \sbox{\lgDiagram}{\input{diagramas/#1}}%
  \ifdim\wd\lgDiagram>\linewidth\resizebox{\linewidth}{!}{\usebox{\lgDiagram}}\else\usebox{\lgDiagram}\fi
  \captionof{figure}{#3}\label{diag:#2}
  \end{minipage}\par\medskip}
\newcommand{\ruta}[1]{\nolinkurl{#1}}
\newcolumntype{L}[1]{>{\raggedright\arraybackslash}p{#1}}
\setlength{\LTpre}{6pt}
\setlength{\LTpost}{9pt}
\setlength{\parskip}{3pt}
\setlength{\emergencystretch}{2em}
\setcounter{tocdepth}{4}
\setcounter{secnumdepth}{3}
\pagestyle{fancy}
\fancyhf{}
\fancyhead[C]{%
  \begin{minipage}[c]{2cm}\centering
      \includegraphics[width=2cm,height=2cm,keepaspectratio]{imagenes/uta.png}
  \end{minipage}\hfill
  \begin{minipage}[c]{8cm}\centering
      \small\textbf{UNIVERSIDAD TÉCNICA DE AMBATO}\\[0.1cm]
      \scriptsize\textit{FACULTAD DE INGENIERÍA EN SISTEMAS, ELECTRÓNICA E INDUSTRIAL}\\[0.1cm]
      \scriptsize\textbf{CARRERA SOFTWARE}\\[0.1cm]
      \scriptsize\textbf{CICLO ACADÉMICO: JULIO--DICIEMBRE 2026}
  \end{minipage}\hfill
  \begin{minipage}[c]{2cm}\centering
      \includegraphics[width=2cm,height=2cm,keepaspectratio]{imagenes/fisei.png}
  \end{minipage}}
\fancyfoot[C]{\small\thepage}
\begin{document}
\begin{center}
    {\large\textbf{INFORME DE TAREA}}
\end{center}

\section{Portada}
\noindent\begin{tabularx}{\linewidth}{@{}L{.37\linewidth} >{\raggedright\arraybackslash}X@{}}
\textbf{Tema:} & LogiMarket 360\\[5pt]
\textbf{Unidad de Organización Curricular:} & PROFESIONAL\\[5pt]
\textbf{Nivel y Paralelo:} & Quinto, paralelo A\\[5pt]
\textbf{Alumnos participantes:} & García Amores Manolo José\\[5pt]
\textbf{Asignatura:} & Patrones de software\\[5pt]
\textbf{Docente:} & Javier Vargas\\
\end{tabularx}
\par\medskip
\noindent\textbf{Práctica:} Guía APE 04 \hfill
\textbf{Modalidad:} Individual

\section{Informe de la Tarea}
\subsection{Objetivos}
\subsubsection{General}
Diseñar, implementar, probar y documentar una aplicación empresarial de escritorio con persistencia real, trazabilidad operativa y arquitectura limpia. Aplicar Clean Code, KISS, DRY, YAGNI, los principios SOLID y los patrones GoF en operaciones verificables.

\subsubsection{Específicos}
\begin{itemize}
\item Modelar ventas y logística a partir de los actores, invariantes, permisos y cifras oficiales de la práctica.
\item Construir cinco proyectos .NET con dependencias verificables y demostrar los 23 patrones dentro de casos de uso.
\item Verificar consistencia ante fallos, recuperación de SQLite, idempotencia, seguridad y resultados mediante pruebas y evidencias.
\end{itemize}

\subsection{Instrucciones}
\begin{itemize}
\item Seguir los requisitos y comprobaciones de la guía APE 04 \cite{ape04}. Ejecutar la aplicación Windows Forms con SQLite, autorización por operación y colaboradores simulados.
\item Usar los datos oficiales y SQLite independiente para cada escenario que exige stock inicial. Reutilizar la base del propio escenario para las pruebas de reinicio, concurrencia o recuperación.
\item Vincular los resultados a código, casos de aceptación, artefactos reproducibles e historial Git. Identificar las condiciones pendientes de validación.
\end{itemize}

\subsection{Listado de equipos, materiales y recursos}
\textbf{Listado de equipos y materiales generales empleados en la tarea:}
\begin{itemize}
\item Equipo Windows, .NET SDK 8, C\#, Windows Forms, EF Core 8, SQLite, xUnit y Git.
\item Guía APE 04, plantilla institucional, repositorio GitHub, registros de pruebas y capturas en una base de demostración aislada.
\item Tectonic para el informe, BibTeX para las referencias y Poppler para la revisión del PDF.
\end{itemize}
\textbf{TAC (Tecnologías para el Aprendizaje y Conocimiento) empleados en la tarea:}
\begin{itemize}
\item[$\square$] Plataformas educativas
\item[$\square$] Simuladores y laboratorios virtuales
\item[$\square$] Aplicaciones educativas
\item[$\square$] Recursos audiovisuales
\item[$\square$] Gamificación
\item[$\checkmark$] Inteligencia Artificial, como apoyo técnico y editorial con contraste de las evidencias.
\item Otros (Especifique): \underline{GitHub y documentación técnica oficial}.
\end{itemize}

\subsection{Desarrollo de la Tarea (Actividades a Desarrollar)}
Se analizaron reglas, flujos y wireframes. Se implementaron el dominio y SQLite, después el núcleo transaccional y los patrones restantes. Se conectaron las doce vistas con Application y se realizó una auditoría de arquitectura, reglas, pruebas y entrega. La comprobación posterior de DPI amplió la cobertura visual. Su alcance está en los resultados.

\clearpage
\begingroup\small\setlength{\parskip}{0pt}\tableofcontents\endgroup
\clearpage
\subsection{Resultados obtenidos}
\subsubsection{Resumen ejecutivo}
LogiMarket 360 controla la venta y entrega de accesorios informáticos mediante un catálogo con kits anidados, pedidos con instantáneas, cobro simulado, existencias y devoluciones. Persiste en SQLite los efectos comerciales y los eventos que deben procesarse después de confirmar. La interfaz utiliza Windows Forms y compara los permisos visibles con la autorización efectiva de Application.

El informe documenta la revisión \ruta{12458bc1cf073e2575e1d85bda0f9f0517a1fa19}. La ejecución automatizada archivada acredita 119 pruebas aprobadas. La matriz de auditoría registra los 38 casos de aceptación y los 23 patrones. El resultado global de entrega sigue condicionado por la verificación de presentación al 125\,\% y al 150\,\%; este requisito es transversal y se detalla en la sección~\ref{sec:dpi}.

\subsubsection{Problema, actores y requisitos}
La empresa necesita controlar existencias físicas al vender tanto hojas como kits, conservar el valor realmente acordado y evitar cobros o reembolsos repetidos. El diseño separa la operación comercial, sus efectos confirmados y la notificación recuperable. Los comprobantes, el proveedor de pago, los avisos y la guía de envío son simulados e internos. La fuente normativa es la guía de práctica \cite{ape04}.

\diagrama{contexto}{contexto}{Contexto operativo en un equipo Windows y sus límites.}
\diagrama{casos}{casos}{Casos de uso y responsabilidades principales por rol.}

\begin{longtable}{@{}L{.19\linewidth}L{.76\linewidth}@{}}
\toprule Rol & Operaciones y alcance \\ \midrule\endhead
ADMIN & Todas las operaciones, incluyendo mantenimiento, demostración y auditoría. \\
VENTAS & Gestión de clientes y pedidos, validación, confirmación y cobro. También consulta el catálogo operativo. \\
BODEGA & Consultas de preparación e inventario, despacho y entrega. El Proxy deniega un intento de cobro antes de las trazas y del proveedor. \\
CONSULTA & Operaciones de lectura, incluidos los reportes permitidos. Sin mutaciones comerciales ni ajustes. \\
\bottomrule
\end{longtable}
\noindent Los permisos se verifican con \ruta{AutorizacionOperaciones.ExigirAsync}. Las pruebas con sesiones por rol contrastan las operaciones de escritura y de consulta. La Figura~\ref{diag:casos} representa el alcance principal, que se completa con las consultas autorizadas en Application.

\paragraph{Trazabilidad funcional y de calidad.}
La siguiente matriz vincula cada requisito con escenarios y evidencia. COMPLETO significa que el alcance auditado cuenta con evidencia. PARCIAL indica una condición pendiente. NO VERIFICADO identifica una ejecución no corroborada. INCUMPLIDO se reserva para una condición que la ejecución contradice.
\small
\begin{longtable}{@{}L{.09\linewidth}L{.47\linewidth}L{.35\linewidth}@{}}
\caption{Trazabilidad de requisitos funcionales y no funcionales.}\\
\toprule ID & Condición de la guía & Evidencia y estado \\ \midrule\endfirsthead
\toprule ID & Condición de la guía & Evidencia y estado \\ \midrule\endhead
RF01 & Identificarse y cerrar sesión localmente con usuarios iniciales, hash y cuatro roles. & P01, P22; login correcto/incorrecto y sesión cerrada. \textbf{COMPLETO} \\[2pt]
RF02 & Consultar y administrar productos con CRUD, búsqueda e invariantes del catálogo. & P01, P21; duplicado SKU y cantidades inválidas. \textbf{COMPLETO} \\[2pt]
RF03 & Registrar y editar kits, impedir ciclos y expandir SKU finales. & P05, P18, P27, P34. \textbf{COMPLETO} \\[2pt]
RF04 & Buscar, registrar y editar clientes con identificación, categoría, contacto y estado activo. & P23; ejercicio CRUD y cliente inactivo. \textbf{COMPLETO} \\[2pt]
RF05 & Crear/editar borrador con cliente, líneas, cantidades, entrega e instantáneas; rechazar vacío. & P08, P13; modificar catálogo posterior y verificar instantánea. \textbf{COMPLETO} \\[2pt]
RF06 & Clonar carrito con copia profunda y deshacer/rehacer hasta cinco estados completos, solo en borrador. & P14, P29. \textbf{COMPLETO} \\[2pt]
RF07 & Calcular subtotal, promoción única, flete, plazo y total con cifras verificables. & P02–P04, P12, P32; umbral 9/10 y redondeo. \textbf{COMPLETO} \\[2pt]
RF08 & Validar por cadena, cancelar únicamente en estados permitidos y consultar historial. & P08, P23, P24, P31. \textbf{COMPLETO} \\[2pt]
RF09 & Confirmar con cobro simulado sustituible, registrar trazas y evitar pago doble/sobreventa. & P02–P07, P16, P19–P20, P22, P36. \textbf{COMPLETO} \\[2pt]
RF10 & Preparar, despachar y entregar en orden, generar guía simulada y consultar seguimiento. & P09, P12, P31; trazabilidad de guía/entrega. \textbf{COMPLETO} \\[2pt]
RF11 & Devolver totalmente desde ENTREGADO conforme a motivo, plazo y condición; reintegro y reembolso únicos. & P33, P35; límites inclusivos y producto usado. \textbf{COMPLETO} \\[2pt]
RF12 & Consultar/registrar inventario y movimientos con motivo y referencia, incluido ajuste auditado. & P02–P07, P35–P36; ajuste positivo/negativo. \textbf{COMPLETO} \\[2pt]
RF13 & Obtener reporte por fechas y filtros por estado, total y cliente con AND/OR y paréntesis. & P25, P26, P34; filtro por rango y resumen. \textbf{COMPLETO} \\[2pt]
RF14 & Emitir comprobantes internos de venta y devolución a vista previa o TXT. & P17; cuatro combinaciones. \textbf{COMPLETO} \\[2pt]
RF15 & Generar avisos simulados EMAIL/SMS, auditoría y métricas por pedido, despacho y devolución. & P11, P15, P19, P28, P30, P37. \textbf{COMPLETO} \\[2pt]
RF16 & Mostrar panel con pedidos, ventas, alertas, stock y métricas; refrescar catálogo, carrito y pedidos mediante mediador. & P15, P21, P28, P30. \textbf{COMPLETO} \\[2pt]
RF17 & Demostrar cambio de proveedor, fallo de cobro y fallo/reintento Outbox en entorno controlado. & P06, P16, P37. \textbf{COMPLETO} \\[2pt]
RF18 & Persistir pedido, pago, stock, despacho, devolución y auditoría al reiniciar. & P01, P10, P35, P37; consulta postreinicio. \textbf{COMPLETO} \\[2pt]
RNF01 & Integridad. Cero efectos de negocio parciales en rechazo. Un pago por pedido, un movimiento por SKU en confirmación, una devolución y stock $\ge$ 0. & P05–P08, P20, P35–P36, índices, estado de DB. \textbf{COMPLETO} \\[2pt]
RNF02 & Mantenibilidad. Cinco proyectos y cero dependencias prohibidas. Patrones, SOLID, Repository, CQRS básico, transacción real y DI por constructor. & P38, \_build\_ limpio, .csproj, diagramas, código e invocaciones reales. \textbf{COMPLETO} \\[2pt]
RNF03 & Usabilidad. Doce áreas accesibles desde la ventana principal; TAB, foco, etiquetas, validación próxima al campo, escala, mensajes, confirmación, selección persistente y dos decimales. & Cotejo de las 12 vistas y alcance en la sección~\ref{sec:dpi}. \textbf{PARCIAL} \\[2pt]
RNF04 & Rendimiento. P21 ejecutado sobre al menos 1.000 filas, con identidad compartida, instancias y memoria frente a escenario equivalente. La guía no fija latencia ni porcentaje de ahorro. & Medición reproducible, condiciones y resultado bruto obtenido. \textbf{COMPLETO} \\[2pt]
RNF05 & Seguridad local. Cuatro roles, autorización de operaciones y cero contraseñas en claro; denegación de pago antes de trazas y proveedor. & P01, P22, inspección de BD, revisión del ZIP. \textbf{COMPLETO} \\[2pt]
RNF06 & Trazabilidad. Registrar cambios de estado, movimientos, correlación, comandos, eventos y evidencias para todos los requisitos/patrones y casos. & Matrices, logs, captura, referencia a archivo/método/commit. \textbf{COMPLETO} \\[2pt]
RNF07 & Testabilidad. Ejecutar P01–P38 y aprobar al menos 25 pruebas automáticas en dominio, aplicación e integración. & P38, informe TRX, datos repetibles y resultados obtenidos. \textbf{COMPLETO} \\[2pt]
RNF08 & Recuperabilidad. Base vacía recuperable vía migraciones/seed, persistencia al reiniciar y evento reintentable sin duplicación. & P01, P10, P37, instrucciones reproducibles. \textbf{COMPLETO} \\[2pt]
RNF09 & Compatibilidad y alcance. Windows, Visual Studio 2022 compatible, C\#/.NET 8, WinForms, EF Core 8, SQLite y \ruta{Microsoft.Extensions.DependencyInjection}. Colaboradores simulados y comprobantes internos. & Compilación, UI, esquema relacional real y demostración sin servicios externos. \textbf{COMPLETO} \\[2pt]
RNF10 & Calidad y autoría. Registrar al menos diez riesgos reales con refactorizaciones, Git y pruebas antes/después. Aplicar cinco principios SOLID y Clean Code/KISS/DRY/YAGNI. & Registro de riesgos e intervenciones en la sección~\ref{sec:riesgos}. \textbf{PARCIAL} \\[2pt]
\bottomrule\end{longtable}\normalsize



\subsubsection{Modelo de dominio y de datos}
\texttt{ProductoSimple} valida SKU, precio, peso y stock. \texttt{KitCompuesto} trabaja con \texttt{IElementoCatalogo} y controla los ciclos al agregar componentes. El precio suma los elementos directos y el desglose expande el grafo hasta productos físicos, con sus multiplicidades. \texttt{Cliente} conserva identidad y categoría NORMAL o VIP. \texttt{Usuario} contiene identidad, rol y hash. MAYORISTA es una política calculada a partir de unidades, independiente del rol y de la categoría persistida.

El agregado \texttt{Pedido} controla sus cambios de estado y guarda detalles, componentes históricos e historial. Las líneas son propias del pedido y preservan identidad de catálogo, precios y cantidades físicas. Los movimientos de inventario guardan el signo, el saldo resultante, el motivo y la referencia. Una devolución puede existir para un pago único. Despacho y entrega se registran como datos operativos con guía y fechas. La auditoría conserva acciones, resultados y correlaciones.

\diagrama{entidades}{entidades}{Clases del catálogo y del pedido con sus relaciones de composición.}
\diagrama{relacional}{relacional}{Relaciones persistidas principales. Se muestran las restricciones que afectan el flujo de venta y recuperación.}
\diagrama{estado}{estado}{Estados y transiciones autorizables del pedido. La devolución registra una operación vinculada tras la entrega.}

\paragraph{Persistencia e integridad.}
\ruta{LogiMarketDbContext} registra los modelos en \ruta{LogiMarket.Infrastructure/Persistencia/Modelo}. La migración \ruta{20261120214119_EsquemaInicial} crea el esquema. \ruta{DatosIniciales} comprueba datos existentes antes de agregar el catálogo de control y los roles. La auditoría archivó el arranque en un archivo nuevo, una verificación de integridad \texttt{ok} y cero errores de claves foráneas.

Los modelos de catálogo, pedidos y operaciones declaran reglas de integridad y unicidad. En operaciones, las claves únicas cubren pedido, idempotencia y referencia de pago, la devolución única y la entrega de un evento por suscriptor. En catálogo, el SKU es único y las restricciones numéricas protegen precio, peso y stock. EF Core configura la persistencia y las transacciones de SQLite \cite{eftransaction,sqlitefk}. La transacción de escritura se abre antes de leer y actualizar los saldos críticos. El aislamiento del escritor en SQLite se comprueba con dos conexiones en P36 \cite{sqliteisolation}.

\subsubsection{Arquitectura y dependencias}
La arquitectura coloca invariantes y reglas del agregado en Domain, orquestación y puertos en Application, adaptadores técnicos en Infrastructure y presentación en WinForms. La inversión de dependencias se resuelve en el punto de composición \cite{cleanarch}. El comportamiento de escritorio se apoya en controles nativos Windows Forms \cite{winforms}.

\diagrama{componentes}{componentes}{Referencias directas entre los cinco proyectos inspeccionados.}
\begin{longtable}{@{}L{.21\linewidth}L{.25\linewidth}L{.47\linewidth}@{}}
\toprule Proyecto & Referencias directas & Responsabilidad y puntos de comprobación \\ \midrule\endhead
Domain & Ninguna propia & \ruta{Pedido}, \ruta{EstadoDelPedido}, catálogo, condiciones del stock e instantáneas. Sin UI ni persistencia. \\
Application & Domain & Casos de uso, puertos y permisos. \ruta{OperacionesPedido} ejecuta comandos. \ruta{ConsultarReportePedidos} y \ruta{ConsultarEscritorio} resuelven consultas. \\
Infrastructure & Application, Domain & EF Core 8, SQLite, migración, repositorios, PBKDF2, simulación de pago, TXT y Outbox. \\
WinForms & Application, Infrastructure & \ruta{Program.CrearHost} compone Infrastructure. \ruta{VentanaComercial} recibe servicios, presenta resultados y navega. \\
Tests & Domain, Application, Infrastructure & xUnit, dobles de puertos, conexiones SQLite reales y pruebas del grafo. \\
\bottomrule
\end{longtable}

\plabel{Repository} \ruta{RepositorioPedidos.BuscarAsync} recupera el agregado con sus detalles; Application depende de \ruta{IRepositorioPedidos}. El catálogo separa \ruta{IConsultaCatalogo} e \ruta{IEscrituraCatalogo}. Las consultas de escritorio y reporte usan proyecciones de sus propios puertos.
\plabel{Unidad de trabajo} \ruta{UnidadDeTrabajoSqlite.EnTransaccionAsync} inicia la transacción, guarda y confirma. Al fallar, el rastreo se limpia y la transacción se dispone sin confirmar. P20 introduce un evento con FK inválida y verifica que pedido, pago, stock y eventos de éxito no queden a medio actualizar.
\plabel{CQRS básico} Los comandos expresan intención comercial y escriben por permisos. Las consultas entregan proyecciones para renderizar. Ambos caminos usan el mismo archivo SQLite; la separación corresponde a responsabilidades, sin añadir un bus externo.
\plabel{Inyección} \ruta{RegistroPersistencia.AgregarPersistenciaLogiMarket} conecta los puertos con los concretos. La composición se verifica construyendo el proveedor con validación de ámbitos y comparando el grafo de los cinco \texttt{.csproj}.


\subsubsection{Flujos de negocio y resultados oficiales}
\paragraph{Acceso, catálogo y clientes.}
El acceso local recupera al usuario, verifica PBKDF2-HMAC-SHA256 y abre sesión con su rol. El hash usa 600\,000 iteraciones, sal de 16 bytes y derivado de 32 bytes. La comparación emplea tiempo fijo. El cierre elimina el actor de sesión. La autorización se aplica antes de las operaciones protegidas, y los intentos denegados se pueden auditar incluso tras una reversión.

\begin{center}\small
\begin{tabular}{@{}llrrr@{}}\toprule
SKU & Producto & USD & Stock inicial & kg\\\midrule
P01 & Mouse USB & 20.00 & 8 & 0.20\\
P02 & Teclado & 30.00 & 5 & 0.70\\
P03 & Memoria USB & 10.00 & 20 & 0.10\\
P04 & Webcam & 40.00 & 4 & 0.35\\
P05 & Alfombrilla & 5.00 & 30 & 0.10\\\bottomrule
\end{tabular}
\captionof{table}{Catálogo oficial desde una base nueva, cotejado en P01.}\label{tab:catalogo-oficial}
\end{center}
KIT\_OFICINA vale USD 50.00. KIT\_TELETRABAJO agrega una P03 a KIT\_OFICINA, vale USD 60.00 y suma 1.00 kg. P18 verifica composiciones y ciclos. P27 y P34 verifican las hojas y el peso.

\paragraph{Pedido, descuentos y entrega.}
El alta construye un borrador con cliente, modalidad, líneas y snapshots. La cadena revisa cliente activo, líneas, existencias y estado. Dos valores de catálogo en momentos distintos pueden coexistir en pedidos diferentes, porque el pedido pagado conserva sus instantáneas. Cada monto monetario usa decimal. Se elige una promoción por prioridad. A partir de diez unidades físicas MAYORISTA aplica 15\,\%. En los demás casos VIP aplica 10\,\% cuando corresponde, o se aplica NORMAL con 0\,\%. El descuento se redondea a dos decimales con \ruta{MidpointRounding.AwayFromZero}. ESTÁNDAR cuesta USD 5.00 y tiene plazo de tres días; EXPRÉS cuesta USD 12.00 y tiene plazo de un día.

\begin{center}\small
\begin{tabular}{@{}L{.22\linewidth}L{.42\linewidth}L{.30\linewidth}@{}}\toprule
Caso & Cálculo en USD & Stock luego del pago\\\midrule
Ana VIP (P02) & $60.00-6.00+5.00=59.00$ & P01=7, P02=4, P03=19\\
Luis NORMAL (P03) & $90.00-0.00+12.00=102.00$ & P04=2, P05=28\\
Carmen MAYORISTA (P04) & $100.00-15.00+5.00=90.00$ & P03=10\\\bottomrule
\end{tabular}
\captionof{table}{Escenarios oficiales ensayados en archivos SQLite separados.}
\end{center}
P32 añade la superposición de VIP y mayorista y prueba el redondeo. Las estrategias son excluyentes.

\paragraph{Confirmación transaccional.}
El invocador registra el comando y la fachada coordina la autorización, la validación y el pago. La cadena entrega las cantidades físicas validadas a la valoración y al descuento de existencias. La pasarela es simulada. Proxy exige permiso antes de Decorator y Adapter. La unidad de trabajo confirma pedido, pago, stock, movimientos y Outbox como una unidad. Un duplicado con igual clave recupera la operación anterior. El rechazado conserva el estado VALIDADO y los saldos. Las trazas diagnósticas distinguen ese rechazo de un cobro registrado.
\diagrama{confirmacion}{confirmacion}{Secuencia aprobada. El procesamiento de Outbox ocurre después de confirmar.}
\diagrama{rechazo}{rechazo}{Secuencia rechazada por falta de stock o por respuesta de la pasarela.}

\paragraph{Entrega, devolución y reporte.}
BODEGA puede despachar un pedido pagado y registrar su entrega. La guía y el plazo se obtienen de la familia elegida. La devolución total requiere ENTREGADO. DEFECTO admite hasta 30 días y ARREPENTIMIENTO hasta 7, en este último caso con producto sin uso. El reintegro se realiza según la decisión de la política. El reembolso reusa el valor histórico descontado sin flete. Ana pagó 59.00, de los cuales 5.00 corresponden a envío, y el reembolso fue de USD 54.00. P35 verifica el reintegro de P01, P02 y P03, la persistencia, el rechazo de un duplicado y la auditoría.

Los reportes consumen un árbol de filtros tipado sobre estado, importe y cliente. El comprobante interno se puede generar como vista previa o TXT. El procesador Outbox reintenta notificaciones, métricas y auditoría a partir del evento confirmado, con deduplicación por suscriptor. P37 comprueba persistencia del intento fallido, reapertura y recuperación sin duplicar el cobro.

\subsubsection{SOLID y programación limpia}
\begin{longtable}{@{}L{.11\linewidth}L{.48\linewidth}L{.35\linewidth}@{}}
\toprule Principio & Evidencia concreta & Comprobación y límite \\ \midrule\endhead
SRP & \ruta{ConfirmacionPedidoFacade.ConfirmarAsync} coordina. \ruta{ValorarPedido.Calcular} resuelve la valoración compuesta. & P02--P07 y P20 distinguen reglas, escritura y rollback. La fachada conserva varios colaboradores. \\
OCP & \ruta{CalculadoraPedido.Calcular} recibe políticas por constructor y elige por prioridad. \ruta{FamiliasEntrega.Obtener} elige la familia registrada. & P32 añade una política y comprueba su resultado. Agregar prioridades iguales requeriría definir el desempate. \\
LSP & \ruta{IResolutorPasarelas.Obtener} entrega \ruta{IPasarelaPago} con adaptador legado o alterno. & P16 ejecuta la misma solicitud en ambos, aprueba 5.21, repite la respuesta y rechaza importes inválidos. Un proveedor externo requeriría conciliación adicional. \\
ISP & \ruta{IConsultaCatalogo}, \ruta{IEscrituraCatalogo}, \ruta{IColaEventos} e \ruta{IProcesadorOutbox} separan consumidores. & P38 verifica firmas y composición. \ruta{ILecturaEscritorio} reúne consultas de apoyo y podría segregarse si crecen sus consumidores. \\
DIP & \ruta{ConfirmacionPedidoFacade} recibe puertos de persistencia, autorización y pasarela; la infraestructura implementa los contratos. & P38 instancia DI y P20 sustituye el puerto de eventos para probar rollback. El registro DI concentra los concretos. \\
\bottomrule
\end{longtable}

Clean Code se concreta en métodos que expresan intención, puertos nombrados por su uso, invariantes encapsuladas y comentarios identificables por patrón. KISS se manifiesta en un único archivo SQLite y separación básica de comandos y consultas. DRY evita repetir la valoración y la expansión de cantidades físicas, comprobado por el cambio \texttt{ccc6cb3}. YAGNI limita el alcance a pago, avisos y guías simuladas, y a comprobantes internos. Los efectos de una operación siguen siendo observables mediante el historial y las pruebas.

\paragraph{Convenciones de los diagramas.}
Un trazo discontinuo con punta abierta indica dependencia; un triángulo hueco, herencia (continuo) o implementación (discontinuo); el rombo lleno marca composición. En secuencias, los retornos y rechazos se distinguen del llamado por trazo discontinuo y, en el rechazo, por color y rótulo.

\clearpage\subsubsection{Patrones GoF creacionales}
Se presenta cada patrón mediante su problema, código real, participantes, prueba y costo. La nomenclatura de las tres familias se toma del catálogo GoF \cite{gof}. Los diagramas muestran las colaboraciones relevantes para el caso de uso.
\paragraph{Factory Method}
\plabel{Problema y propósito} Seleccionar EMAIL o SMS para un evento durable sin acoplar el suscriptor a la construcción de cada emisor.
\plabel{Participantes y colaboración} Application/Eventos/\ruta{AvisosSimulados.cs} \ruta{CreadorEmail.Crear}, \ruta{CreadorSms.Crear} y \ruta{ObservadorAvisos.RecibirAsync}. CreadorEmail y CreadorSms especializan Crear. El observador selecciona por canal, genera el contenido del evento y lo registra mediante el puerto. El aviso se dispara al reprocesar Outbox, después del commit.
\plabel{Intervención en el caso de uso} Outbox selecciona emisor mediante creadores DI.
\diagrama{patron-factory-method}{factory-method}{Colaboración Factory Method en LogiMarket 360.}
\plabel{Código fuente} \ruta{src/LogiMarket.Application/Eventos/AvisosSimulados.cs}, líneas 22--34 del commit auditado.
\lstinputlisting[firstnumber=22]{codigo/factory-method.cs}
\plabel{Prueba y resultado} P11 y P30 comprueban tipo, canal y contenido distintivo. El aviso es un registro simulado, sin correo ni SMS reales. \textbf{COMPLETO}.
\plabel{Evidencia} \ruta{OutboxComercialTests.P11_P30}.
\plabel{Beneficio y costo} Canal intercambiable; dos creadores.
\medskip
\clearpage
\paragraph{Abstract Factory}
\plabel{Problema y propósito} Obtener tarifa, plazo y etiqueta de la misma modalidad y detectar combinaciones inconsistentes.
\plabel{Participantes y colaboración} Application/Entregas/\ruta{FamiliasEntrega.cs} \ruta{FabricaEntrega.Crear}, variantes, \ruta{FamiliasEntrega.Obtener}. FabricaEntregaEstandar y FabricaEntregaExpres crean TarifaEntrega, PlazoEntrega y EtiquetaEntrega. ValorarPedido y DespacharPedidos usan la familia resuelta para cotizar y elaborar la guía.
\plabel{Intervención en el caso de uso} Valorar y despachar reciben componentes de la misma familia.
\diagrama{patron-abstract-factory}{abstract-factory}{Colaboración Abstract Factory en LogiMarket 360.}
\plabel{Código fuente} \ruta{src/LogiMarket.Application/Entregas/FamiliasEntrega.cs}, líneas 29--37 del commit auditado.
\lstinputlisting[firstnumber=29]{codigo/abstract-factory.cs}
\plabel{Prueba y resultado} P12 cotejó USD 5.00 y 3 días para ESTÁNDAR, USD 12.00 y 1 día para EXPRÉS, y guías distinguibles. \textbf{COMPLETO}.
\plabel{Evidencia} \ruta{OperacionesComercialesTests.P12}.
\plabel{Beneficio y costo} Evita mezclar variantes; tres productos por familia.
\medskip
\clearpage
\paragraph{Builder}
\plabel{Problema y propósito} Construir un pedido con identidad, cliente y varias líneas preservando la coherencia del agregado.
\plabel{Participantes y colaboración} Application/Pedidos/\ruta{PedidoBuilder.cs} \ruta{Construir}, \ruta{AdministrarPedidos.CrearAsync} y \ruta{Pedido}. AdministrarPedidos prepara instantáneas y usa el builder para componer las líneas. Construir materializa Pedido con el usuario y la fecha, antes de solicitar persistencia.
\plabel{Intervención en el caso de uso} Construye líneas y snapshots para altas válidas.
\diagrama{patron-builder}{builder}{Colaboración Builder en LogiMarket 360.}
\plabel{Código fuente} \ruta{src/LogiMarket.Application/Pedidos/PedidoBuilder.cs}, líneas 19--30 del commit auditado.
\lstinputlisting[firstnumber=19]{codigo/builder.cs}
\plabel{Prueba y resultado} P13 rechaza pedido vacío y cliente inexistente. La prueba de alta con varios renglones obtiene subtotal USD 67.00. \textbf{COMPLETO}.
\plabel{Evidencia} \ruta{OperacionesComercialesTests.ABMCatalogoClientesYBorradoresPreservaInstantaneas}.
\plabel{Beneficio y costo} Composición explícita; coordina invariantes.
\medskip
\clearpage
\paragraph{Prototype}
\plabel{Problema y propósito} Duplicar un borrador comercial sin trasladar su identidad persistente ni compartir líneas mutables.
\plabel{Participantes y colaboración} Domain/Pedidos/\ruta{Pedido.Borrador.cs} \ruta{ClonarComoBorrador}; Application/Pedidos/\ruta{AdministrarPedidos.cs} \ruta{ClonarBorradorAsync}. ClonarComoBorrador exige edición permitida y crea un pedido nuevo. Las líneas se reconstruyen con copias de sus snapshots. El caso de uso carga y persiste la copia mediante el repositorio.
\plabel{Intervención en el caso de uso} Pedido y líneas se copian con identidad nueva.
\diagrama{patron-prototype}{prototype}{Colaboración Prototype en LogiMarket 360.}
\plabel{Código fuente} \ruta{src/LogiMarket.Domain/Pedidos/Pedido.Borrador.cs}, líneas 8--20 del commit auditado.
\lstinputlisting[firstnumber=8]{codigo/prototype.cs}
\plabel{Prueba y resultado} P14 modifica una copia sin cambiar la plantilla ni otra copia. La evidencia visual conserva identificadores distintos. \textbf{COMPLETO}.
\plabel{Evidencia} \ruta{PatronesBorradorTests.P14}; fase 4, captura 16.
\plabel{Beneficio y costo} Reutilización; copia proporcional al grafo.
\medskip
\clearpage
\paragraph{Singleton}
\plabel{Problema y propósito} Compartir un contador coherente entre los observadores y dos vistas, aunque se resuelvan en distintos ámbitos de DI.
\plabel{Participantes y colaboración} Application/Eventos/\ruta{MetricasSesion.cs} \ruta{Registrar}, \ruta{Contar}; Infrastructure/Persistencia/\ruta{RegistroPersistencia.cs}. El registro de persistencia instancia MetricasSesion como singleton del host. El acumulador deduplica por evento. ConsultarEscritorio proyecta métricas en Panel y Auditoría.
\plabel{Intervención en el caso de uso} DI comparte una instancia entre Outbox y consultas de ambas vistas.
\diagrama{patron-singleton}{singleton}{Colaboración Singleton en LogiMarket 360.}
\plabel{Código fuente} \ruta{src/LogiMarket.Application/Eventos/MetricasSesion.cs}, líneas 6--18 del commit auditado.
\lstinputlisting[firstnumber=6]{codigo/singleton.cs}
\plabel{Prueba y resultado} P15 y P37 verifican identidad, 50 eventos únicos en paralelo y el valor 1 en dos vistas de una misma sesión. Tras cerrar el proceso el contador de sesión empieza de nuevo. \textbf{COMPLETO}.
\plabel{Evidencia} \ruta{MetricasSesionTests.P15}; \ruta{OutboxComercialTests.P37}; fase 4, capturas 06–07.
\plabel{Beneficio y costo} Coherencia por proceso; memoria por evento.
\medskip
\clearpage
\subsubsection{Patrones GoF estructurales}
\paragraph{Adapter}
\plabel{Problema y propósito} Aprovechar el simulador legado que recibe centavos y entregar una respuesta uniforme al caso de uso.
\plabel{Participantes y colaboración} Infrastructure/Pagos/\ruta{PasarelasSimuladas.cs} \ruta{PasarelaLegadaAdapter.CobrarAsync}, \ruta{ResolutorPasarelas.Obtener}. El adaptador convierte decimal a unidades menores, valida precisión a través del conversor y normaliza la respuesta. ResolutorPasarelas también puede seleccionar PasarelaAlternaSimulada con el mismo contrato.
\plabel{Intervención en el caso de uso} Adapta llamadas de confirmación y admite alterno.
\diagrama{patron-adapter}{adapter}{Colaboración Adapter en LogiMarket 360.}
\plabel{Código fuente} \ruta{src/LogiMarket.Infrastructure/Pagos/PasarelasSimuladas.cs}, líneas 25--35 del commit auditado.
\lstinputlisting[firstnumber=25]{codigo/adapter.cs}
\plabel{Prueba y resultado} P16 prueba ambos proveedores, valor USD 5.21, repetición y rechazo de importes inválidos. \textbf{COMPLETO}.
\plabel{Evidencia} \ruta{ConfirmacionTests.P16}.
\plabel{Beneficio y costo} Sustitución; valida precisión al convertir.
\medskip
\clearpage
\paragraph{Bridge}
\plabel{Problema y propósito} Combinar dos tipos de documento interno con dos medios sin duplicar una clase por cada combinación.
\plabel{Participantes y colaboración} Application/Documentos/\ruta{Comprobante.cs} \ruta{EmitirAsync}; \ruta{EmitirComprobantes.EjecutarAsync}; Infrastructure/Documentos/\ruta{ArchivoTXT.cs}. ComprobanteVenta y ComprobanteDevolucion preparan DocumentoInterno. Comprobante.EmitirAsync delega el contenido al puerto ISalidaDocumento, implementado por VistaPrevia y ArchivoTXT.
\plabel{Intervención en el caso de uso} Dos tipos de comprobante $\times$ vista y archivo.
\diagrama{patron-bridge}{bridge}{Colaboración Bridge en LogiMarket 360.}
\plabel{Código fuente} \ruta{src/LogiMarket.Application/Documentos/Comprobante.cs}, líneas 26--35 del commit auditado.
\lstinputlisting[firstnumber=26]{codigo/bridge.cs}
\plabel{Prueba y resultado} P17 cubre cuatro combinaciones. El paquete conserva comprobantes TXT reales de USD 59.00 y USD 54.00. \textbf{COMPLETO}.
\plabel{Evidencia} \ruta{ComprobantesTests.BridgeEntregaVentaYDevolucionACadaUnoDeLosDosMedios}; TXT y registro visual de fase 3.
\plabel{Beneficio y costo} Cambios independientes; contrato de documento común.
\medskip
\clearpage
\paragraph{Composite}
\plabel{Problema y propósito} Tratar productos y kits como elementos uniformes, incluyendo kits dentro de kits.
\plabel{Participantes y colaboración} Domain/Catalogo/\ruta{IElementoCatalogo.cs}, \ruta{ProductoSimple}, \ruta{KitCompuesto.CalcularPrecio}, \ruta{AgregarComponente}. ProductoSimple resuelve el precio de una hoja. KitCompuesto suma precios de sus elementos con cantidades, comprueba ciclos al agregar y genera snapshots. La confirmación descuenta SKU físicos.
\plabel{Intervención en el caso de uso} Seed compone kits y fachada descuenta hojas.
\diagrama{patron-composite}{composite}{Colaboración Composite en LogiMarket 360.}
\plabel{Código fuente} \ruta{src/LogiMarket.Domain/Catalogo/KitCompuesto.cs}, líneas 79--89 del commit auditado.
\lstinputlisting[firstnumber=79]{codigo/composite.cs}
\plabel{Prueba y resultado} P18 confirma precios 50.00 y 60.00, anidación y rechazo de ciclos. P05 comprueba que el stock del componente P02 gobierna la confirmación. \textbf{COMPLETO}.
\plabel{Evidencia} \ruta{InicializacionTests.P01}; CatalogoTests; RestriccionesSqliteTests.
\plabel{Beneficio y costo} Composición anidada; necesita controlar ciclos.
\medskip
\clearpage
\paragraph{Decorator}
\plabel{Problema y propósito} Correlacionar el intento de pago, la llamada al proveedor y el desenlace sin cambiar su contrato.
\plabel{Participantes y colaboración} Application/Pagos/\ruta{AuditoriaPagoDecorator.cs} \ruta{CobrarAsync}, \ruta{ITrazasPago}. AuditoriaPagoDecorator escribe etapas mediante ITrazasPago y delega al proveedor envuelto. La fachada guarda la traza fuera de la transacción comercial concluida o revertida.
\plabel{Intervención en el caso de uso} Envuelve al proveedor y correlaciona etapas.
\diagrama{patron-decorator}{decorator}{Colaboración Decorator en LogiMarket 360.}
\plabel{Código fuente} \ruta{src/LogiMarket.Application/Pagos/AuditoriaPagoDecorator.cs}, líneas 7--20 del commit auditado.
\lstinputlisting[firstnumber=7]{codigo/decorator.cs}
\plabel{Prueba y resultado} P19 compara INICIO, PASARELA, FIN o ERROR. ITrazasPago es intercambiable y queda el riesgo residual de persistencia diagnóstica R12. \textbf{COMPLETO}.
\plabel{Evidencia} \ruta{ConfirmacionTests.P06_P19} y P16.
\plabel{Beneficio y costo} Auditoría transversal; persistencia posterior al commit.
\medskip
\clearpage
\paragraph{Facade}
\plabel{Problema y propósito} Exponer una operación única para confirmar, aunque intervengan permisos, existencias, pago y persistencia.
\plabel{Participantes y colaboración} Application/Pagos/\ruta{ConfirmacionPedidoFacade.cs} \ruta{ConfirmarAsync} y \ruta{ConfirmarEnTransaccionAsync}. ConfirmarAsync exige permiso y abre unidad de trabajo. ConfirmarEnTransaccionAsync resuelve idempotencia, revisa el pedido y consulta la pasarela. Una aprobación aplica el pago, modifica saldos y encola Outbox.
\plabel{Intervención en el caso de uso} Valida y valora, registra pago, stock y Outbox.
\diagrama{patron-facade}{facade}{Colaboración Facade en LogiMarket 360.}
\plabel{Código fuente} \ruta{src/LogiMarket.Application/Pagos/ConfirmacionPedidoFacade.cs}, líneas 50--65 del commit auditado.
\lstinputlisting[firstnumber=50]{codigo/facade.cs}
\plabel{Prueba y resultado} P02–P07 comprueban cálculo, rechazo e idempotencia. P20 prueba reversión cuando falla la persistencia previa al commit. \textbf{COMPLETO}.
\plabel{Evidencia} \ruta{ConfirmacionTests.P02_P03_P04}. \ruta{ConfirmacionTests.P05}. \ruta{ConfirmacionTests.P06_P19}. \ruta{ConfirmacionTests.P07_P08}. \ruta{ConfirmacionTests.P20}.
\plabel{Beneficio y costo} Punto de entrada estable; varios colaboradores.
\medskip
\clearpage
\paragraph{Flyweight}
\plabel{Problema y propósito} Reducir objetos repetidos de nombre, categoría y descripción cuando un mismo SKU se representa muchas veces.
\plabel{Participantes y colaboración} Application/Catalogo/\ruta{CatalogoVisualFlyweightFactory.cs} \ruta{CrearRenglon}, \ruta{ConsultarCatalogoVisual.DemostrarAsync}. Cada renglón mantiene precio, cantidad y stock propios. La fábrica interna comparte un objeto inmutable por SKU. La comparación usó la misma cantidad y orden de filas y metadatos equivalentes.
\plabel{Intervención en el caso de uso} Comparte metadatos e independiza precio, cantidad y saldo.
\diagrama{patron-flyweight}{flyweight}{Colaboración Flyweight en LogiMarket 360.}
\plabel{Código fuente} \ruta{src/LogiMarket.Application/Catalogo/CatalogoVisualFlyweightFactory.cs}, líneas 18--31 del commit auditado.
\lstinputlisting[firstnumber=18]{codigo/flyweight.cs}
\plabel{Prueba y resultado} P21 verificó identidad de referencia y 5 frente a 1.000 metadatos. La ejecución archivada midió 56.752 frente a 104.024 bytes asignados. El valor corresponde a esa medición, no a todo el proceso. \textbf{COMPLETO}.
\plabel{Evidencia} \ruta{FlyweightTests.P21}; fase 4, captura 10; protocolo fase 3.
\plabel{Beneficio y costo} Menos asignaciones con repetición; diccionario por proyección.
\medskip
\clearpage
\paragraph{Proxy}
\plabel{Problema y propósito} Evitar que roles de consulta o bodega inicien un cobro o induzcan trazas de la pasarela.
\plabel{Participantes y colaboración} Application/Pagos/\ruta{PagoAutorizadoProxy.cs} \ruta{CobrarAsync}; Infrastructure/Pagos/\ruta{PasarelasSimuladas.cs} \ruta{Obtener}. El resolutor arma Proxy, Decorator y Adapter/alterno en ese orden. Proxy vuelve a verificar Operacion.Cobrar inmediatamente antes de acceder al decorador.
\plabel{Intervención en el caso de uso} Exige autorización antes del Decorator y Adapter.
\diagrama{patron-proxy}{proxy}{Colaboración Proxy en LogiMarket 360.}
\plabel{Código fuente} \ruta{src/LogiMarket.Application/Pagos/PagoAutorizadoProxy.cs}, líneas 8--13 del commit auditado.
\lstinputlisting[firstnumber=8]{codigo/proxy.cs}
\plabel{Prueba y resultado} P22 verifica cero llamadas y cero etapas de cobro para CONSULTA y BODEGA, con auditoría de seguridad. \textbf{COMPLETO}.
\plabel{Evidencia} \ruta{ConfirmacionTests.P22}.
\plabel{Beneficio y costo} Bloquea efectos; requiere orden correcto de DI.
\medskip
\clearpage
\subsubsection{Patrones GoF comportamentales}
\paragraph{Chain of Responsibility}
\plabel{Problema y propósito} Detener la confirmación en el primer requisito de negocio aplicable y conservar un código diagnóstico.
\plabel{Participantes y colaboración} Application/Pedidos/\ruta{CadenaValidacionPedido.cs} \ruta{RevisarAsync}, \ruta{Paso.EjecutarAsync}. La fachada exige el rol. La cadena comprueba cliente, líneas, stock y estado. El resultado incluye las cantidades físicas validadas, que se reutilizan en valoración e inventario.
\plabel{Intervención en el caso de uso} Cliente $\to$ líneas $\to$ stock $\to$ estado; permiso en fachada.
\diagrama{patron-chain-of-responsibility}{chain-of-responsibility}{Colaboración Chain of Responsibility en LogiMarket 360.}
\plabel{Código fuente} \ruta{src/LogiMarket.Application/Pedidos/CadenaValidacionPedido.cs}, líneas 21--33 del commit auditado.
\lstinputlisting[firstnumber=21]{codigo/chain-of-responsibility.cs}
\plabel{Prueba y resultado} P23 verifica CLIENTE antes que STOCK y ausencia de efectos de negocio y de pasarela. La prueba examina la fase de precondiciones. \textbf{COMPLETO}.
\plabel{Evidencia} \ruta{OperacionesComercialesTests.P23}, validación y fachada. Cotejo del cliente ausente: \ruta{cotejo-p23.trx}.
\plabel{Beneficio y costo} Reglas aisladas; orden semántico.
\medskip
\clearpage
\paragraph{Command}
\plabel{Problema y propósito} Representar confirmar, cancelar, despachar o devolver como una acción que la interfaz y las pruebas puedan invocar de la misma forma.
\plabel{Participantes y colaboración} Application/Comandos/\ruta{ComandosPedido.cs}, \ruta{OperacionesPedido.ConfirmarAsync}, \ruta{InvocadorComandos.EjecutarAsync}. OperacionesPedido liga cada comando con su caso receptor. El invocador ejecuta, conserva correlación y registra el resultado. La interfaz despacha el mismo tipo de comando que las pruebas.
\plabel{Intervención en el caso de uso} Confirmar, cancelar, despachar y devolver quedan trazables.
\diagrama{patron-command}{command}{Colaboración Command en LogiMarket 360.}
\plabel{Código fuente} \ruta{src/LogiMarket.Application/Comandos/InvocadorComandos.cs}, líneas 16--32 del commit auditado.
\lstinputlisting[firstnumber=16]{codigo/command.cs}
\plabel{Prueba y resultado} P24 usa a Ana y el proveedor LEGADO. El resultado USD 59.00 se asocia a ConfirmarPedidoCommand y a auditoría de éxito. \textbf{COMPLETO}.
\plabel{Evidencia} \ruta{OutboxComercialTests.P24}; recorrido P24 en evidencia de fase 3 y TXT de venta.
\plabel{Beneficio y costo} Separa invocador de receptor; auditoría posterior.
\medskip
\clearpage
\paragraph{Interpreter}
\plabel{Problema y propósito} Permitir expresiones del usuario sin convertir su texto en SQL.
\plabel{Participantes y colaboración} Application/Reportes/\ruta{ParserFiltroPedido.cs} \ruta{Parsear}, \ruta{FiltroPedidos.Cumple}, \ruta{ConsultarReportePedidos.EjecutarAsync}. El parser construye nodos de estado, total y cliente, combina AND antes de OR y respeta paréntesis. ConsultarReportePedidos evalúa el árbol mediante Cumple. Los errores señalan la columna.
\plabel{Intervención en el caso de uso} Parser produce árbol y reporte lo aplica.
\diagrama{patron-interpreter}{interpreter}{Colaboración Interpreter en LogiMarket 360.}
\plabel{Código fuente} \ruta{src/LogiMarket.Application/Reportes/ParserFiltroPedido.cs}, líneas 15--28 del commit auditado.
\lstinputlisting[firstnumber=15]{codigo/interpreter.cs}
\plabel{Prueba y resultado} P25 recupera solo el pedido coincidente. P26 verifica sintaxis incorrecta, precedencia, paréntesis y error localizable. \textbf{COMPLETO}.
\plabel{Evidencia} \ruta{FiltroPedidosTests.P25}; recorrido de fase 4. \ruta{FiltroPedidosTests.P26}; cierre fase 4, captura 20.
\plabel{Beneficio y costo} Precedencia verificable; lectura materializada.
\medskip
\clearpage
\paragraph{Iterator}
\plabel{Problema y propósito} Recorrer los componentes físicos de kits anidados sin ceder la colección mutable ni perder multiplicidades.
\plabel{Participantes y colaboración} Domain/Catalogo/\ruta{IteradorComponentesFisicos.cs} \ruta{Recorrer}; \ruta{AdministrarCatalogo.ConsultarExistenciasAsync}. Recorrer mantiene el estado de recorrido y entrega hojas en orden determinista, propagando factores de cantidad. La consulta de existencias por kit usa ese resultado.
\plabel{Intervención en el caso de uso} Enumera hojas con multiplicidad en consulta real.
\diagrama{patron-iterator}{iterator}{Colaboración Iterator en LogiMarket 360.}
\plabel{Código fuente} \ruta{src/LogiMarket.Domain/Catalogo/IteradorComponentesFisicos.cs}, líneas 8--23 del commit auditado.
\lstinputlisting[firstnumber=8]{codigo/iterator.cs}
\plabel{Prueba y resultado} P27 devuelve P01, P02 y P03 en orden. También se comprobaron ramas compartidas y multiplicidades sin bucles. \textbf{COMPLETO}.
\plabel{Evidencia} \ruta{RecorridosCatalogoTests.P27_P34} e IteratorEsDeterminista.
\plabel{Beneficio y costo} Oculta estado de recorrido; pila e instantánea.
\medskip
\clearpage
\paragraph{Mediator}
\plabel{Problema y propósito} Actualizar vistas dependientes sin que una pantalla conozca a las demás.
\plabel{Participantes y colaboración} Application/Navegacion/\ruta{FormularioPrincipalMediator.cs} \ruta{NotificarAsync}/\ruta{Destinos}; WinForms/\ruta{VentanaComercial.cs} \ruta{NotificarAsync}. Destinos decide qué secciones dependen del cambio y NotificarAsync ejecuta las suscripciones. El flujo de confirmación sincroniza catálogo, carrito, pedidos, inventario y métricas.
\plabel{Intervención en el caso de uso} Panel, auditoría, inventario, carrito y pedidos se refrescan mediante suscripciones.
\diagrama{patron-mediator}{mediator}{Colaboración Mediator en LogiMarket 360.}
\plabel{Código fuente} \ruta{src/LogiMarket.Application/Navegacion/FormularioPrincipalMediator.cs}, líneas 36--48 del commit auditado.
\lstinputlisting[firstnumber=36]{codigo/mediator.cs}
\plabel{Prueba y resultado} P28 verifica cinco colaboradores. La regresión de despacho probó que el contador bajara de 1 a 0 y que Outbox pasara de 1 a 2 después del arreglo. \textbf{COMPLETO}.
\plabel{Evidencia} \ruta{MediadorComercialTests.ConfirmacionRefrescaCincoColaboradoresSobreElEstadoConfirmado}; recorrido de fase 3.
\plabel{Beneficio y costo} Sin llamadas entre vistas; mantener destinos al añadir eventos.
\medskip
\clearpage
\paragraph{Memento}
\plabel{Problema y propósito} Deshacer y rehacer borradores sin permitir revertir un pago ni modificar el historial de estados.
\plabel{Participantes y colaboración} Domain/Pedidos/\ruta{Pedido.Borrador.cs} \ruta{CapturarBorrador}/\ruta{RestaurarBorrador}, Application/Pedidos/\ruta{HistorialBorrador.cs}. Pedido es el originador y copia sus líneas en un memento. HistorialBorrador guarda hasta cinco estados y controla versiones. La restauración se autoriza solo mientras admite edición.
\plabel{Intervención en el caso de uso} Historial acotado que deshace y rehace con control de versión.
\diagrama{patron-memento}{memento}{Colaboración Memento en LogiMarket 360.}
\plabel{Código fuente} \ruta{src/LogiMarket.Domain/Pedidos/Pedido.Borrador.cs}, líneas 22--36 del commit auditado.
\lstinputlisting[firstnumber=22]{codigo/memento.cs}
\plabel{Prueba y resultado} P29 comprueba tres ediciones, dos retrocesos, un avance, límite de cinco y bloqueo tras confirmar. \textbf{COMPLETO}.
\plabel{Evidencia} \ruta{PatronesBorradorTests.P29}; fase 4, captura 17.
\plabel{Beneficio y costo} Reversión aislada; memoria por pedido activo.
\medskip
\clearpage
\paragraph{Observer}
\plabel{Problema y propósito} Notificar distintos efectos de un hecho de negocio confirmado, con recuperación por suscriptor.
\plabel{Participantes y colaboración} Application/Eventos/\ruta{PublicadorEventos.cs} \ruta{PublicarAsync}; Infrastructure/Persistencia/Eventos/\ruta{EntregaDurableSqlite.cs} \ruta{EntregarAsync}. PublicadorEventos recorre a los observadores de avisos, métricas y auditoría. EntregaDurableSqlite separa la entrega por suscriptor y preserva marcas persistentes según el tipo de efecto.
\plabel{Intervención en el caso de uso} Outbox despacha a auditoría, avisos y métricas.
\diagrama{patron-observer}{observer}{Colaboración Observer en LogiMarket 360.}
\plabel{Código fuente} \ruta{src/LogiMarket.Application/Eventos/PublicadorEventos.cs}, líneas 31--36 del commit auditado.
\lstinputlisting[firstnumber=31]{codigo/observer.cs}
\plabel{Prueba y resultado} P30 confirma los tres efectos y ausencia de evento exitoso tras rollback. P37 confirma recuperación y deduplicación. \textbf{COMPLETO}.
\plabel{Evidencia} \ruta{OutboxComercialTests.P11_P30} y P37; \ruta{ConfirmacionTests.P20}. \ruta{OutboxComercialTests.P37}; fase 4, capturas 02/06.
\plabel{Beneficio y costo} Suscriptores ampliables; marcas idempotentes.
\medskip
\clearpage
\paragraph{State}
\plabel{Problema y propósito} Encapsular qué acciones admite el pedido en cada etapa y registrar solo transiciones autorizadas.
\plabel{Participantes y colaboración} Domain/Pedidos/\ruta{EstadoDelPedido.cs} \ruta{Para}, Domain/Pedidos/\ruta{Pedido.cs}. EstadoDelPedido.Para resuelve el comportamiento del valor persistible. Borrador admite editar, validar o cancelar. Validado admite pagar o cancelar. Pagado y Despachado avanzan; los estados cerrados deniegan los cambios.
\plabel{Intervención en el caso de uso} Agregado delega acciones en comportamiento del estado.
\diagrama{patron-state}{state}{Colaboración State en LogiMarket 360.}
\plabel{Código fuente} \ruta{src/LogiMarket.Domain/Pedidos/EstadoDelPedido.cs}, líneas 29--45 del commit auditado.
\lstinputlisting[firstnumber=29]{codigo/state.cs}
\plabel{Prueba y resultado} P31 recorre el ciclo permitido y rechaza repetición, retroceso y cancelación después de pagar. \textbf{COMPLETO}.
\plabel{Evidencia} \ruta{PedidoTests.StateRegistraSoloTransicionesAutorizadasPorElEstado}.
\plabel{Beneficio y costo} Centraliza transiciones; mapa persistencia/comportamiento.
\medskip
\clearpage
\paragraph{Strategy}
\plabel{Problema y propósito} Seleccionar una política vigente en función de prioridad, cliente y cantidad física final.
\plabel{Participantes y colaboración} Application/Precios/\ruta{CalculadoraPedido.cs} \ruta{Calcular}, \ruta{IPoliticaDescuento}. Las estrategias SinDescuento, Vip10 y Mayorista15 reciben el contexto. Calcular filtra aplicables, elige una por prioridad y redondea el descuento con AwayFromZero.
\plabel{Intervención en el caso de uso} Valorar elige mayorista, VIP o normal sin acumular.
\diagrama{patron-strategy}{strategy}{Colaboración Strategy en LogiMarket 360.}
\plabel{Código fuente} \ruta{src/LogiMarket.Application/Precios/CalculadoraPedido.cs}, líneas 46--55 del commit auditado.
\lstinputlisting[firstnumber=46]{codigo/strategy.cs}
\plabel{Prueba y resultado} P32 obtiene 0.00, 6.00 y 15.00 en los escenarios correspondientes. Prueba que VIP y MAYORISTA no se acumulan y cubre una política nueva. \textbf{COMPLETO}.
\plabel{Evidencia} \ruta{OperacionesComercialesTests.P32} y PromocionYUmbralFisico.
\plabel{Beneficio y costo} Extensión por DI; mantener prioridades inequívocas.
\medskip
\clearpage
\paragraph{Template Method}
\plabel{Problema y propósito} Conservar la misma secuencia de devolución para varios motivos, con plazos y condiciones propios.
\plabel{Participantes y colaboración} Application/Devoluciones/\ruta{ProcesarDevolucion.cs} \ruta{EjecutarAsync} y variantes. El método plantilla ordena elegibilidad, inspección, cálculo histórico, registro, reintegro y evento. Las variantes DevolucionPorDefecto y DevolucionPorArrepentimiento resuelven la condición particular.
\plabel{Intervención en el caso de uso} Misma secuencia para elegibilidad, reembolso, stock y evento.
\diagrama{patron-template-method}{template-method}{Colaboración Template Method en LogiMarket 360.}
\plabel{Código fuente} \ruta{src/LogiMarket.Application/Devoluciones/ProcesarDevolucion.cs}, líneas 44--60 del commit auditado.
\lstinputlisting[firstnumber=44]{codigo/template-method.cs}
\plabel{Prueba y resultado} P33 acepta DEFECTO día 20 y rechaza ARREPENTIMIENTO día 10 sin efectos. P35 comprueba USD 54.00 sin flete y devolución única. \textbf{COMPLETO}.
\plabel{Evidencia} \ruta{LogisticaTests.P33}. \ruta{LogisticaTests.P35}; capturas usuario-inventario, usuario-devolucion y usuario-auditoria.
\plabel{Beneficio y costo} Mantiene orden; limitar pasos variables.
\medskip
\clearpage
\paragraph{Visitor}
\plabel{Problema y propósito} Añadir operaciones de desglose y peso a árboles de catálogo sin integrarlas en las entidades.
\plabel{Participantes y colaboración} Domain/Catalogo/\ruta{VisitantesCatalogo.cs}, \ruta{ProductoSimple.Aceptar}, \ruta{KitCompuesto.Aceptar}; \ruta{ConsultarLogisticaAsync}. Aceptar distingue hoja y kit. VisitanteComponentes propaga cantidades. VisitanteDesgloseSku acumula unidades y VisitantePesoLogistico suma el peso físico en la consulta logística.
\plabel{Intervención en el caso de uso} El catálogo aplica visitantes de desglose y peso.
\diagrama{patron-visitor}{visitor}{Colaboración Visitor en LogiMarket 360.}
\plabel{Código fuente} \ruta{src/LogiMarket.Domain/Catalogo/VisitantesCatalogo.cs}, líneas 12--26 del commit auditado.
\lstinputlisting[firstnumber=12]{codigo/visitor.cs}
\plabel{Prueba y resultado} P34 obtiene P01, P02 y P03 y peso 1.00 kg. La suite comprueba también multiplicidad por tres. \textbf{COMPLETO}.
\plabel{Evidencia} \ruta{RecorridosCatalogoTests.P27_P34}; fase 4, captura 09.
\plabel{Beneficio y costo} Reportes extensibles; nuevo nodo afecta el contrato.
\medskip



\clearpage
\subsubsection{Interfaz Windows Forms y comprobaciones funcionales}
Las doce vistas siguen las operaciones de los wireframes de las páginas 11--16 de la guía. Los formularios muestran datos de Application y protegen visualmente las acciones. La autorización decisiva se verifica en el caso de uso. Las capturas corresponden a ejecuciones reales y se identifican por recorrido. La revisión incorporó siete recapturas a 96 DPI sobre copias aisladas de los escenarios. El registro gráfico identifica las tomas históricas que aún conservan el puntero o su halo. El manifiesto de evidencias relaciona cada archivo con su origen. Algunas se tomaron en fases previas al último ajuste DPI; se usan para documentar la operación observada.

\paragraph{1. Acceso}
\plabel{Objetivo} Volver a identificar al operador o finalizar la sesión local.
\plabel{Controles} Usuario, contraseña enmascarada, acción de acceso y mensaje junto al campo.
\plabel{Permisos} Antes de autenticar no hay rol operativo. Cada una de las cuatro cuentas se vincula a su rol.
\plabel{Servicios} \ruta{AutenticacionLocal.IniciarAsync}; \ruta{SesionLocal}.
\plabel{Patrones y decisiones} Autenticación y autorización por operación.
\plabel{Comprobación} Se observó el acceso centrado y la vuelta al formulario al cerrar sesión. El usuario confirmó ingreso con las cuatro cuentas. La captura es del formulario de producción, sin credenciales expuestas. Véase la Figura~\ref{fig:acceso}.
\vistafig{acceso-revision.jpg}{acceso}{10.1cm}{Formulario de acceso a 96 DPI.}

\clearpage
\paragraph{2. Panel general}
\plabel{Objetivo} Recuperar el estado de pedidos, ventas, stock bajo y eventos pendientes.
\plabel{Controles} Tarjetas, tabla de pedidos recientes y acceso al pedido seleccionado.
\plabel{Permisos} Cualquier sesión habilitada puede consultar.
\plabel{Servicios} \ruta{ConsultarEscritorio.ResumenAsync}.
\plabel{Patrones y decisiones} Singleton, Mediator y Observer.
\plabel{Comprobación} Se observó Ana, el estado y el total de USD 59.00 después de reabrir. La comparación de métricas entre Panel y Auditoría se detalla en P15. Véase la Figura~\ref{fig:panel}.
\vistafig{panel-revision.jpg}{panel}{10.1cm}{Panel general con datos del escenario comercial de demostración.}

\clearpage
\paragraph{3. Catálogo y kits}
\plabel{Objetivo} Buscar SKU y mantener productos, kits y sus componentes.
\plabel{Controles} Buscador, tablas, edición, desglose del kit y acción para medir 1.000 filas.
\plabel{Permisos} Consulta para el rol autenticado. Mantenimiento reservado a ADMIN.
\plabel{Servicios} \ruta{AdministrarCatalogo}; \ruta{ConsultarCatalogoVisual.DemostrarAsync}; \ruta{ConsultarEscritorio.KitsAsync}.
\plabel{Patrones y decisiones} Composite, Iterator, Visitor y Flyweight.
\plabel{Comprobación} El kit anidado devolvió P01/P02/P03 y 1.00 kg. Se modificó el precio de P01 y se comparó con el pedido ya pagado. La acción de edición quedó inhabilitada con VENTAS. Véase la Figura~\ref{fig:catalogo}.
\vistafig{catalogo-cierre.jpg}{catalogo}{10.1cm}{Desglose real de KIT\_TELETRABAJO y sus componentes físicos.}

\clearpage
\paragraph{4. Gestión de clientes}
\plabel{Objetivo} Buscar, registrar y editar clientes por identificación o nombre.
\plabel{Controles} Buscador, tabla, ficha de cliente y acciones Nuevo/Editar.
\plabel{Permisos} ADMIN y VENTAS gestionan clientes. Los otros roles quedan sujetos al permiso de lectura.
\plabel{Servicios} \ruta{AdministrarClientes.BuscarAsync}; \ruta{AdministrarClientes}.
\plabel{Patrones y decisiones} Segregación entre consulta y escritura y autorización en Application.
\plabel{Comprobación} Se observó navegación de Buscar cliente al botón Buscar y ficha de Ana VIP. La prueba de operaciones comerciales cubre la escritura y el uso del cliente en el pedido. Véase la Figura~\ref{fig:clientes}.
\vistafig{clientes-revision.jpg}{clientes}{10.1cm}{Búsqueda y datos de cliente con sesión VENTAS.}

\clearpage
\paragraph{5. Carrito y borradores}
\plabel{Objetivo} Crear pedidos, editar líneas y conservar snapshots.
\plabel{Controles} Selectores de pedido, cliente, entrega y SKU; cantidad; edición, clonado, deshacer, rehacer y validación.
\plabel{Permisos} ADMIN y VENTAS pueden editar y validar. El agregado impide los cambios fuera de BORRADOR.
\plabel{Servicios} \ruta{AdministrarPedidos}; \ruta{AdministrarCatalogo}; \ruta{ConsultarEscritorio.DetalleAsync}.
\plabel{Patrones y decisiones} Builder, Prototype, Memento, Composite y Mediator.
\plabel{Comprobación} La copia recibió una identidad nueva. La cantidad cambió 1 a 2 y luego volvió a 1 y a 2 mediante deshacer y rehacer. El cliente y los precios del original se conservaron. Véase la Figura~\ref{fig:carrito}.
\vistafig{carrito-cierre.jpg}{carrito}{10.1cm}{Restauración de líneas en un borrador durante el recorrido de Memento.}

\clearpage
\paragraph{6. Pedidos e historial}
\plabel{Objetivo} Consultar pedidos, estados y precio histórico, y cancelar cuando corresponda.
\plabel{Controles} Filtro por estado, tabla de pedidos, detalle, historial y acción Cancelar.
\plabel{Permisos} Consulta por roles autenticados. ADMIN y VENTAS cancelan en un estado permitido.
\plabel{Servicios} \ruta{ConsultarReportePedidos.EjecutarAsync}; \ruta{ConsultarEscritorio.DetalleAsync}; \ruta{OperacionesPedido.CancelarAsync}.
\plabel{Patrones y decisiones} State, Interpreter, Command y Mediator.
\plabel{Comprobación} La captura documenta protección del pedido pagado. P08 verifica denegación de edición/cancelación y preservación del importe, pago e historial tras cambiar el precio de catálogo. Véase la Figura~\ref{fig:pedidos}.
\vistafig{pedidos.jpg}{pedidos}{10.1cm}{Pedido pagado y protección del histórico.}

\clearpage
\paragraph{7. Confirmación y cobro}
\plabel{Objetivo} Mostrar cifras, seleccionar un simulador y confirmar el pago de un pedido válido.
\plabel{Controles} Tabla de órdenes, selector de proveedor, resumen y diálogo de confirmación.
\plabel{Permisos} ADMIN y VENTAS pueden cobrar. BODEGA y CONSULTA reciben denegación en Application.
\plabel{Servicios} \ruta{ConsultarEscritorio.DetalleAsync}; \ruta{OperacionesPedido.ConfirmarAsync}.
\plabel{Patrones y decisiones} Command, Facade, Chain, Proxy, Decorator, Adapter, Strategy y Abstract Factory.
\plabel{Comprobación} Con Ana se ejecutó el cobro simulado por USD 59.00. Con BODEGA la acción quedó inactiva y legible. P22 comprueba la denegación efectiva y ausencia de llamadas al proveedor. Véase la Figura~\ref{fig:cobro}.
\vistafig{cobro-revision.jpg}{cobro}{10.1cm}{Resumen de cobro visible con la acción de confirmación inhabilitada para BODEGA.}

\clearpage
\paragraph{8. Despacho y entrega}
\plabel{Objetivo} Emitir una guía, registrar entrega y consultar seguimiento.
\plabel{Controles} Selector de pedido, guía, plazo, historial y botones de despacho/entrega.
\plabel{Permisos} ADMIN y BODEGA escriben despacho/entrega. VENTAS y CONSULTA solo acceden a las consultas permitidas.
\plabel{Servicios} \ruta{DespacharPedidos.ConsultarAsync}; \ruta{OperacionesPedido.DespacharAsync}; \ruta{DespacharPedidos.EntregarAsync}.
\plabel{Patrones y decisiones} State, Abstract Factory, Command y Mediator.
\plabel{Comprobación} Se registró una guía ESTÁNDAR de tres días y se observó el estado ENTREGADO. P09 prueba rechazo antes del pago y no autorización de VENTAS. Véase la Figura~\ref{fig:despacho}.
\vistafig{despacho-revision.jpg}{despacho}{10.1cm}{Guía emitida y seguimiento del pedido entregado.}

\clearpage
\paragraph{9. Devoluciones}
\plabel{Objetivo} Consultar elegibilidad y registrar el reembolso total sin flete.
\plabel{Controles} Selector de motivo, casilla de producto sin uso, Validar política, Confirmar devolución y tarjeta de importe.
\plabel{Permisos} La devolución solo se registra con ADMIN. Application verifica el rol y la política.
\plabel{Servicios} \ruta{AdministrarDevoluciones.PrevisualizarAsync}; \ruta{OperacionesPedido.DevolverAsync}.
\plabel{Patrones y decisiones} Template Method, Command, State y Observer vía Outbox.
\plabel{Comprobación} La captura aportada por el usuario muestra devolución registrada tras reabrir, cobro USD 59.00 y reembolso USD 54.00. Se observaron controles inactivos ante una segunda devolución. P35 comprueba unicidad también en SQLite. Véase la Figura~\ref{fig:devoluciones}.
\vistafig{devoluciones.png}{devoluciones}{10.1cm}{Devolución de Ana leída tras reiniciar. El importe excluye USD 5.00 de flete.}

\clearpage
\paragraph{10. Inventario y movimientos}
\plabel{Objetivo} Consultar stock y trazabilidad de movimientos por SKU.
\plabel{Controles} Buscador, tabla de stock, tabla de movimientos y registrar ajuste.
\plabel{Permisos} Roles autenticados consultan. ADMIN registra ajustes.
\plabel{Servicios} \ruta{AdministrarCatalogo.AjustarAsync}; \ruta{ConsultarEscritorio.MovimientosAsync}.
\plabel{Patrones y decisiones} Iterator en consultas de kits y Mediator en actualización de la vista.
\plabel{Comprobación} Los tres SKU de Ana mostraron movimientos de venta −1 y devolución +1. Al finalizar, P01/P02/P03 recuperaron 8/5/20. CONSULTA vio la operación de ajuste inactiva en el cierre de fase. Véase la Figura~\ref{fig:inventario}.
\vistafig{inventario.png}{inventario}{10.1cm}{Captura del usuario con los saldos reintegrados y los seis movimientos relacionados.}

\clearpage
\paragraph{11. Reportes y comprobantes}
\plabel{Objetivo} Consultar pedidos por período o expresión y emitir un documento interno.
\plabel{Controles} Rango temporal, filtro de expresiones, tabla, tipo, vista previa y Guardar TXT.
\plabel{Permisos} Consulta y emisión interna sujetas al permiso Consultar.
\plabel{Servicios} \ruta{ConsultarReportePedidos}; \ruta{EmitirComprobantes.EjecutarAsync}; \ruta{IFabricaArchivoTxt}.
\plabel{Patrones y decisiones} Interpreter y Bridge.
\plabel{Comprobación} Un filtro por estado PAGADO y total mínimo recuperó a Ana. Se generó vista previa de USD 59.00 y un TXT real. La prueba P17 cubre venta y devolución en ambos medios. Véase la Figura~\ref{fig:reportes}.
\vistafig{reportes-cierre.jpg}{reportes}{10.1cm}{Vista previa de un comprobante interno de venta.}

\clearpage
\paragraph{12. Auditoría, métricas y demostración}
\plabel{Objetivo} Revisar operaciones, indicadores de sesión y eventos pendientes.
\plabel{Controles} Tarjetas, filtro de categoría, tablas de auditoría y Outbox, selector de canal y controles de simulación/reproceso.
\plabel{Permisos} Auditoría y demostración reservadas a ADMIN. Los fallos inducidos requieren modo demostración.
\plabel{Servicios} \ruta{ConsultarEscritorio}; \ruta{ControlDemostracion.AplicarAsync}; \ruta{ReprocesarEventos.EjecutarAsync}.
\plabel{Patrones y decisiones} Observer, Factory Method, Singleton y Mediator.
\plabel{Comprobación} La captura del usuario registra confirmación y devolución aprobadas en la misma referencia. P37 y las capturas complementarias documentan un fallo y su recuperación en un segundo intento. Véase la Figura~\ref{fig:auditoria}.
\vistafig{auditoria.png}{auditoria}{10.1cm}{Auditoría de comandos de pago y devolución de Ana en demostración.}

\clearpage
\paragraph{Cotejo de DPI, foco y accesibilidad}\label{sec:dpi}
A 96 DPI se verificaron acceso, panel y auditoría, incluida la distribución maximizada y foco representativo. Al 125\,\% se observó \ruta{DeviceDpi=120} en la versión auditada y se abrieron las doce vistas en ventana normal y maximizada. Los dos rótulos largos se leen completos. Se observó el avance por TAB de Buscar cliente a Buscar y el borde de foco. Las vistas largas se alcanzan mediante desplazamiento, aunque algunas celdas largas usan elipsis en ventana normal.

\vistafig{dpi125-foco.jpg}{dpi-foco}{9.1cm}{Foco de teclado visible en Buscar, con la interfaz a 120 DPI.}
La Figura~\ref{fig:dpi-foco} documenta el control activo. Este recorrido representa la navegación en esa vista. Las recapturas editoriales a 96 DPI no amplían esta comprobación de 120 DPI.

\clearpage
Los JPEG archivados a 120 DPI muestran las cuatro cifras de Auditoría en ventana normal, maximizada y al volver a la vista. La bitácora de ejecución registró una omisión que no se reproduce al examinar los píxeles de esos archivos. Falta confirmar físicamente si el efecto aparece durante la navegación. El 150\,\% (144 DPI) tras la corrección sigue NO VERIFICADO. Las capturas del error previo muestran el motivo de la corrección, pero se requiere evidencia posterior.

\vistafig{dpi125-reentrada.jpg}{dpi-normal}{10.3cm}{Reentrada en Auditoría a 120 DPI. El JPEG conserva los valores 0, 0, 0 y 3.}
\clearpage
\vistafig{dpi125-max.jpg}{dpi-max}{9.62cm}{Auditoría maximizada a 120 DPI, con los cuatro valores visibles.}
Las Figuras~\ref{fig:dpi-normal} y \ref{fig:dpi-max} documentan las capturas. RNF03 sigue PARCIAL. Véase \ruta{evidencias/cotejo-dpi.md}.



\clearpage
\subsubsection{Pruebas automatizadas y evidencia de ejecución}
El proyecto \ruta{tests/LogiMarket.Tests/LogiMarket.Tests.csproj} usa xUnit 2.9.3 y Microsoft.NET.Test.Sdk 17.8.0. La suite se inicia con \ruta{dotnet test} y recoge los resultados por caso en TRX \cite{xunit}. Las pruebas de integración usan SQLite real con archivos temporales. Cada entorno se construye con dependencias registradas en el proyecto de persistencia y usa la misma base solo cuando el caso requiere repetición, concurrencia o reapertura.

\begin{center}
\begin{tabular}{@{}lrrrr@{}}\toprule
Grupo & Ejecutadas & Aprobadas & Fallidas & Omitidas\\\midrule
Aplicacion & 1 & 1 & 0 & 0 \\
Arquitectura & 7 & 7 & 0 & 0 \\
Dominio & 21 & 21 & 0 & 0 \\
Integracion & 90 & 90 & 0 & 0 \\
\midrule Total & 119 & 119 & 0 & 0\\\bottomrule
\end{tabular}
\captionof{table}{Conteo de definiciones por espacio de nombres y resumen de \texttt{auditoria-final.trx}.}
\end{center}

\begin{longtable}{@{}L{.22\linewidth}L{.73\linewidth}@{}}
\toprule Área & Clases y comprobaciones \\ \midrule\endhead
Dominio & \ruta{CatalogoTests} y \ruta{PedidoTests}. Precios, stock, SKU, snapshots, ciclos y transiciones. \\
Aplicación & \ruta{MetricasSesionTests}. Identidad y deduplicación concurrente del contador. \\
Integración SQLite & \ruta{InicializacionTests}, \ruta{RestriccionesSqliteTests} y \ruta{PersistenciaTests}. Migración, semilla repetible, datos recuperados y restricciones efectivas. \\
Seguridad & \ruta{SeguridadComercialTests} y \ruta{RevocacionCobroTests}. Sesión, permisos efectivos por rol y guardado de la denegación tras rollback. \\
Transacción & \ruta{ConfirmacionTests}. Escenarios oficiales, proveedor alterno, clave repetida, rollback por FK y competencia por una unidad en dos conexiones. \\
Logística y eventos & \ruta{LogisticaTests} y \ruta{OutboxComercialTests}. Despacho, devolución, stock, observadores, reintento y reapertura. \\
Interfaz y regresiones & \ruta{EscritorioTests}, \ruta{MediadorComercialTests}, \ruta{PatronesBorradorTests}, \ruta{FiltroPedidosTests} y \ruta{ComprobantesTests}. Integración de servicios subyacentes a las pantallas. La observación de la interfaz se registra aparte. \\
Arquitectura & \ruta{ReferenciasTests} y \ruta{ContratosEscritorioTests}. Cinco proyectos, referencias, puertos segregados y composición. \\
\bottomrule
\end{longtable}

La auditoría archivó un build Release con cero errores y cero advertencias. El conteo de pruebas procede del TRX, cuyo resumen y huella se pueden comprobar en el expediente. \ruta{auditoria-base.trx} tiene 118 aprobadas. \ruta{mediator-antes.trx} y \ruta{mediator-despues.trx} prueban la corrección de despacho. \ruta{expansion-antes.trx} y \ruta{expansion-despues.trx} conservan 19 aprobadas en la selección afectada por la refactorización. Se trata de ejecuciones históricas; la compilación LaTeX no las sustituye.

Se volvió a ejecutar la solución del commit auditado sobre una copia limpia durante la preparación de este informe. Con SDK 8.0.428 se obtuvieron cinco proyectos restaurados, compilación Release con cero errores y cero advertencias y 119 de 119 pruebas aprobadas. Se conservan \ruta{auditoria-revalidada.trx} y las salidas de restauración, compilación y prueba. La restauración usó paquetes existentes y desactivó NuGet Audit, por lo que ese comando no aporta una revisión de vulnerabilidades de dependencias.

El cotejo adicional de P23 ejecutó la cadena con cliente ausente sobre una base de prueba y luego con el cliente restaurado. Produjo CLIENTE y después STOCK. La única prueba focalizada aprobó y se entrega con su código y \ruta{cotejo-p23.trx}; es una ejecución separada y queda fuera de las 119 de la suite oficial.

\paragraph{Reproducción en Windows.}
Ejecutar desde la raíz de la solución con .NET 8 y el componente de escritorio:
\begin{verbatim}
dotnet restore LogiMarket360.sln
dotnet build LogiMarket360.sln -c Release --no-restore
dotnet test LogiMarket360.sln -c Release --no-build
\end{verbatim}
Para conservar un TRX, agregar \ruta{--logger} con valor \ruta{trx;LogFileName=logimarket360.trx}. Los artefactos de arranque documentan la migración, los recuentos del seed y la integridad referencial sin adjuntar una base operativa.

\paragraph{Cobertura consolidada de GoF}
\begin{longtable}{@{}L{.25\linewidth}L{.18\linewidth}L{.50\linewidth}@{}}
\toprule Familia & Cobertura & Escenarios y evidencia \\ \midrule\endhead
Creacionales & 5/5 & P11--P15, composición por DI y métricas en dos vistas. \\
Estructurales & 7/7 & P16--P22, simuladores, comprobantes internos, persistencia y medición Flyweight. \\
Comportamentales & 11/11 & P23--P34, relojes y bases aisladas, recuperación de eventos y actualización de vistas. \\
\midrule Total & 23/23 & Participación en casos de uso y comprobaciones individuales documentadas.\\
\bottomrule
\end{longtable}




\paragraph{Definición del alcance.}
El TRX permite comprobar ejecución y resultado sin atribuir a una prueba observable la calidad del foco o la disposición en DPI. Las capturas muestran un estado concreto del proceso, mientras que P36 y P37 demuestran propiedades de concurrencia y recuperación utilizando los datos y distancias temporales de sus propios casos.



\subsubsection{Aceptación P01--P38}
Cada caso conserva el procedimiento previsto, el resultado observado y la evidencia del corte auditado. Las ejecuciones históricas se identifican con su fase. Las evidencias visuales y su alcance se describen por separado. P01--P38 sustenta el alcance de aceptación de la guía; la verificación transversal de DPI se evalúa en RNF03.
\subsection*{Funcionalidad y reglas}
\par\medskip\noindent\begin{minipage}{\linewidth}
\noindent\textbf{\large P01}\quad\textsf{COMPLETO}
\plabel{Descripción} Migración y seed repetible.
\plabel{Precondiciones y entrada} Base ausente; ejecutar seed dos veces.
\plabel{Resultado esperado} Migraciones aplicadas; cinco productos, dos kits y roles; sin duplicados al relanzar seed.
\plabel{Resultado obtenido} Esquema vigente, 5 productos, 2 kits, 4 roles y sin duplicados.
\plabel{Método y evidencia} \ruta{InicializacionTests.P01_SeedDesdeArchivoAusenteEsIdempotenteYPersiste}. \ruta{InicializacionTests.P01}; sqlite-arranque-\ruta{limpio.json}.
\end{minipage}\par
\par\medskip\noindent\begin{minipage}{\linewidth}
\noindent\textbf{\large P02}\quad\textsf{COMPLETO}
\plabel{Descripción} Compra de Ana VIP.
\plabel{Precondiciones y entrada} Ana VIP, teletrabajo $\times$1 estándar, base nueva.
\plabel{Resultado esperado} Total 59.00; stock P01=7, P02=4, P03=19; estado PAGADO y un pago.
\plabel{Resultado obtenido} 59.00, stocks 7/4/19, PAGADO y pago único.
\plabel{Método y evidencia} \ruta{ConfirmacionTests.P02_P03_P04_EscenariosOficiales}. \ruta{ConfirmacionTests.P02_P03_P04}.
\end{minipage}\par
\par\medskip\noindent\begin{minipage}{\linewidth}
\noindent\textbf{\large P03}\quad\textsf{COMPLETO}
\plabel{Descripción} Compra de Luis NORMAL.
\plabel{Precondiciones y entrada} Luis NORMAL, P04 $\times$2 y P05 $\times$2 exprés, base nueva.
\plabel{Resultado esperado} Total 102.00; P04=2, P05=28; un pago.
\plabel{Resultado obtenido} 102.00, stocks P04/P05 2/28, pago único.
\plabel{Método y evidencia} \ruta{ConfirmacionTests.P02_P03_P04_EscenariosOficiales}. \ruta{ConfirmacionTests.P02_P03_P04}.
\end{minipage}\par
\par\medskip\noindent\begin{minipage}{\linewidth}
\noindent\textbf{\large P04}\quad\textsf{COMPLETO}
\plabel{Descripción} Compra mayorista de Carmen.
\plabel{Precondiciones y entrada} Carmen, P03 $\times$10 estándar, base nueva.
\plabel{Resultado esperado} Total 90.00; P03=10; descuento 15.00.
\plabel{Resultado obtenido} 90.00, descuento 15.00, P03=10, pago único.
\plabel{Método y evidencia} \ruta{ConfirmacionTests.P02_P03_P04_EscenariosOficiales}. \ruta{ConfirmacionTests.P02_P03_P04}.
\end{minipage}\par
\par\medskip\noindent\begin{minipage}{\linewidth}
\noindent\textbf{\large P05}\quad\textsf{COMPLETO}
\plabel{Descripción} Stock insuficiente.
\plabel{Precondiciones y entrada} Teletrabajo $\times$6, stock oficial.
\plabel{Resultado esperado} Rechazo por P02 insuficiente; sin pago/stock/evento de éxito.
\plabel{Resultado obtenido} P02 insuficiente; sin pago, movimientos ni Outbox.
\plabel{Método y evidencia} \ruta{ConfirmacionTests.P05_KitSinStockNoInvocaProveedor}. \ruta{ConfirmacionTests.P05}.
\end{minipage}\par
\par\medskip\noindent\begin{minipage}{\linewidth}
\noindent\textbf{\large P06}\quad\textsf{COMPLETO}
\plabel{Descripción} Cobro rechazado.
\plabel{Precondiciones y entrada} Pedido válido y alterno configurado para rechazar.
\plabel{Resultado esperado} Se conserva inventario y estado previo; sin pago aprobado ni Outbox de éxito.
\plabel{Resultado obtenido} VALIDADO y saldos intactos; sin pago ni evento aprobado.
\plabel{Método y evidencia} \ruta{ConfirmacionTests.P06_P19_RechazoMantieneValidadoYRegistraTrazaEnOrden}. \ruta{ConfirmacionTests.P06_P19}.
\end{minipage}\par
\par\medskip\noindent\begin{minipage}{\linewidth}
\noindent\textbf{\large P07}\quad\textsf{COMPLETO}
\plabel{Descripción} Repetición de confirmación.
\plabel{Precondiciones y entrada} Reintentar misma clave en el mismo pedido.
\plabel{Resultado esperado} Un pago y un movimiento por SKU; devuelve el resultado original sin cobro doble.
\plabel{Resultado obtenido} Mismo pago, una llamada y tres movimientos físicos.
\plabel{Método y evidencia} \ruta{ConfirmacionTests.P07_P08_IdempotenciaYProteccionDelPedidoPagado}. \ruta{ConfirmacionTests.P07_P08}.
\end{minipage}\par
\par\medskip\noindent\begin{minipage}{\linewidth}
\noindent\textbf{\large P08}\quad\textsf{COMPLETO}
\plabel{Descripción} Inmutabilidad después del pago.
\plabel{Precondiciones y entrada} Cambiar precio del catálogo y editar/cancelar pedido pagado.
\plabel{Resultado esperado} Ambas operaciones denegadas; datos e historial intactos.
\plabel{Resultado obtenido} Denegación de ambas operaciones y mismo precio, pago e historial.
\plabel{Método y evidencia} TRX focalizado y capturas originales 13–15. \ruta{EscritorioTests.P08}; capturas fase 4 13–15.
\end{minipage}\par
\par\medskip\noindent\begin{minipage}{\linewidth}
\noindent\textbf{\large P09}\quad\textsf{COMPLETO}
\plabel{Descripción} Despacho y entrega por rol.
\plabel{Precondiciones y entrada} BODEGA: intentar antes del pago; pagar, despachar y entregar.
\plabel{Resultado esperado} Éxito solo en secuencia permitida; rechazo sin cambios en estado inválido.
\plabel{Resultado obtenido} Rechazo previo; estado, guía y entrega correctos; VENTAS no despacha.
\plabel{Método y evidencia} \ruta{LogisticaTests.P09_BodegaDespachaEntregaYConsultaGuiaSoloTrasPago}. \ruta{LogisticaTests.P09}; cierre fase 4, captura 17.
\end{minipage}\par
\par\medskip\noindent\begin{minipage}{\linewidth}
\noindent\textbf{\large P10}\quad\textsf{COMPLETO}
\plabel{Descripción} Persistencia después de reiniciar.
\plabel{Precondiciones y entrada} Cerrar y reabrir sobre misma SQLite.
\plabel{Resultado esperado} Información persiste en SQLite y coincide con operaciones anteriores.
\plabel{Resultado obtenido} Pedido, pago, stock y traza preservados; comprobación posterior al reembolso.
\plabel{Método y evidencia} Capturas originales 04–06, captura de cierre 11 y consulta de solo lectura 02. Fase 4 04–06 y SQLite 03/08; usuario-devolucion-\ruta{reabierta.png}.
\end{minipage}\par
\subsection*{Patrones creacionales}
\par\medskip\noindent\begin{minipage}{\linewidth}
\noindent\textbf{\large P11}\quad\textsf{COMPLETO}
\plabel{Descripción} Factory Method por canal.
\plabel{Precondiciones y entrada} Procesar evento por EMAIL y SMS en bases separadas.
\plabel{Resultado esperado} Cada creador retorna notificador concreto y registra el aviso correcto.
\plabel{Resultado obtenido} Registro del canal, tipo y contenido característico.
\plabel{Método y evidencia} \ruta{OutboxComercialTests.P11_P30_CanalElegidoYTresObservadoresRecibenUnEvento}. \ruta{OutboxComercialTests.P11_P30}.
\end{minipage}\par
\par\medskip\noindent\begin{minipage}{\linewidth}
\noindent\textbf{\large P12}\quad\textsf{COMPLETO}
\plabel{Descripción} Abstract Factory de entregas.
\plabel{Precondiciones y entrada} Resolver estándar y exprés.
\plabel{Resultado esperado} Tarifas 5/12 y plazos 3/1 días; guías coherentes con la familia.
\plabel{Resultado obtenido} 5/12, 3/1 y guías de su propia familia.
\plabel{Método y evidencia} \ruta{OperacionesComercialesTests.P12_FabricaProduceFamiliaCoherente}. \ruta{OperacionesComercialesTests.P12}.
\end{minipage}\par
\par\medskip\noindent\begin{minipage}{\linewidth}
\noindent\textbf{\large P13}\quad\textsf{COMPLETO}
\plabel{Descripción} Builder de pedido.
\plabel{Precondiciones y entrada} Crear pedido sin líneas, con cliente inexistente y de varios renglones.
\plabel{Resultado esperado} Incompleto rechazado; completo crea agregado válido.
\plabel{Resultado obtenido} Dos rechazos y agregado válido (subtotal 67.00).
\plabel{Método y evidencia} \ruta{OperacionesComercialesTests.ABMCatalogoClientesYBorradoresPreservaInstantaneas}. \ruta{OperacionesComercialesTests.ABMCatalogoClientesYBorradoresPreservaInstantaneas}.
\end{minipage}\par
\par\medskip\noindent\begin{minipage}{\linewidth}
\noindent\textbf{\large P14}\quad\textsf{COMPLETO}
\plabel{Descripción} Prototype de borrador.
\plabel{Precondiciones y entrada} Clonar borrador con kit anidado dos veces, editar una copia.
\plabel{Resultado esperado} Copia profunda; plantilla/copia previa permanecen intactas.
\plabel{Resultado obtenido} Identidades nuevas y snapshot aislado de plantilla y otra copia.
\plabel{Método y evidencia} \ruta{PatronesBorradorTests.P14_ClonaInstantaneaDeKitSinIdentidadesYPermiteEdicionIndependiente}. \ruta{PatronesBorradorTests.P14}; fase 4, captura 16.
\end{minipage}\par
\par\medskip\noindent\begin{minipage}{\linewidth}
\noindent\textbf{\large P15}\quad\textsf{COMPLETO}
\plabel{Descripción} Singleton de métricas.
\plabel{Precondiciones y entrada} Dos ámbitos y dos vistas; duplicar 50 eventos en paralelo.
\plabel{Resultado esperado} Misma instancia compartida, contadores observables y thread-safe.
\plabel{Resultado obtenido} Identidad única; total 50; mismo contador visible en Panel y Auditoría.
\plabel{Método y evidencia} TRX focalizado y capturas originales 06–07. \ruta{MetricasSesionTests.P15}; \ruta{OutboxComercialTests.P37}; fase 4, capturas 06–07.
\end{minipage}\par
\subsection*{Patrones estructurales}
\par\medskip\noindent\begin{minipage}{\linewidth}
\noindent\textbf{\large P16}\quad\textsf{COMPLETO}
\plabel{Descripción} Adapter y sustitución del proveedor.
\plabel{Precondiciones y entrada} Resolver legado y alterno con el mismo contrato.
\plabel{Resultado esperado} Cobros desde ambos con mismo contrato; legado intacto.
\plabel{Resultado obtenido} Ambos aprueban 5.21, repiten resultado y rechazan importes inválidos.
\plabel{Método y evidencia} \ruta{ConfirmacionTests.P16_AmbosProveedoresCompartenContrato}. \ruta{ConfirmacionTests.P16}.
\end{minipage}\par
\par\medskip\noindent\begin{minipage}{\linewidth}
\noindent\textbf{\large P17}\quad\textsf{COMPLETO}
\plabel{Descripción} Bridge de comprobantes.
\plabel{Precondiciones y entrada} Emitir venta y devolución por VistaPrevia y ArchivoTXT.
\plabel{Resultado esperado} Cuatro combinaciones; VistaPrevia visible y TXT real.
\plabel{Resultado obtenido} Cuatro combinaciones; contenido de venta 59.00 y devolución 54.00.
\plabel{Método y evidencia} \ruta{ComprobantesTests.BridgeEntregaVentaYDevolucionACadaUnoDeLosDosMedios}. \ruta{ComprobantesTests.BridgeEntregaVentaYDevolucionACadaUnoDeLosDosMedios}; TXT y registro visual de fase 3.
\end{minipage}\par
\par\medskip\noindent\begin{minipage}{\linewidth}
\noindent\textbf{\large P18}\quad\textsf{COMPLETO}
\plabel{Descripción} Composite de catálogo.
\plabel{Precondiciones y entrada} Kits oficiales y ciclos directo e indirecto.
\plabel{Resultado esperado} Precios 50/60; expansión correcta; ciclo rechazado.
\plabel{Resultado obtenido} 50/60, hojas P01/P02/P03, rechazo de ciclos.
\plabel{Método y evidencia} \ruta{CatalogoTests}, \ruta{InicializacionTests}, \ruta{RestriccionesSqliteTests}. \ruta{InicializacionTests.P01}; CatalogoTests; RestriccionesSqliteTests.
\end{minipage}\par
\par\medskip\noindent\begin{minipage}{\linewidth}
\noindent\textbf{\large P19}\quad\textsf{COMPLETO}
\plabel{Descripción} Decorator y trazas de cobro.
\plabel{Precondiciones y entrada} Cobro autorizado aprobado, rechazado y argumento inválido.
\plabel{Resultado esperado} LOG de inicio y fin/error rodea la ejecución, no altera resultado.
\plabel{Resultado obtenido} Orden INICIO $\to$ PASARELA $\to$ FIN o ERROR; correlación y resultados.
\plabel{Método y evidencia} \ruta{ConfirmacionTests.P06_P19_RechazoMantieneValidadoYRegistraTrazaEnOrden; P16_AmbosProveedoresCompartenContrato}. \ruta{ConfirmacionTests.P06_P19} y P16.
\end{minipage}\par
\par\medskip\noindent\begin{minipage}{\linewidth}
\noindent\textbf{\large P20}\quad\textsf{COMPLETO}
\plabel{Descripción} Facade y fallo de persistencia.
\plabel{Precondiciones y entrada} Forzar FK inválida en Outbox antes del commit.
\plabel{Resultado esperado} Operación consistente; ninguna persistencia parcial indebida.
\plabel{Resultado obtenido} Pedido VALIDADO; pago, stock, movimientos y evento sin efectos parciales.
\plabel{Método y evidencia} \ruta{ConfirmacionTests.P20_FalloDePersistenciaReviertePagoEstadoMovimientosYEvento}. \ruta{ConfirmacionTests.P20}.
\end{minipage}\par
\par\medskip\noindent\begin{minipage}{\linewidth}
\noindent\textbf{\large P21}\quad\textsf{COMPLETO}
\plabel{Descripción} Flyweight de metadatos.
\plabel{Precondiciones y entrada} 1.000 filas con cinco SKU, mismo orden y datos extrínsecos.
\plabel{Resultado esperado} ReferenceEquals true; identidad y memoria documentadas.
\plabel{Resultado obtenido} ReferenceEquals, 5/1.000 objetos y 56.752/104.024 bytes asignados en medición original.
\plabel{Método y evidencia} \ruta{FlyweightTests.P21_MilRenglonesCompartenCincoInstanciasYMidenAsignaciones}. \ruta{FlyweightTests.P21}; fase 4, captura 10; protocolo fase 3.
\end{minipage}\par
\par\medskip\noindent\begin{minipage}{\linewidth}
\noindent\textbf{\large P22}\quad\textsf{COMPLETO}
\plabel{Descripción} Proxy y roles restringidos.
\plabel{Precondiciones y entrada} CONSULTA y BODEGA intentan cobrar.
\plabel{Resultado esperado} Rechazo antes de Decorator/Adapter; sin cobro ni LOG de cobro.
\plabel{Resultado obtenido} Denegación anterior a log de pago y pasarela; solo auditoría de seguridad.
\plabel{Método y evidencia} \ruta{ConfirmacionTests.P22_ProxyDeniegaAntesDelDecoradorYDeLaPasarela}. \ruta{ConfirmacionTests.P22}.
\end{minipage}\par
\subsection*{Patrones comportamentales}
\par\medskip\noindent\begin{minipage}{\linewidth}
\noindent\textbf{\large P23}\quad\textsf{COMPLETO}
\plabel{Descripción} Chain y orden de validación.
\plabel{Precondiciones y entrada} Kit sin stock con cliente inactivo y, en el cotejo focalizado, cliente ausente al evaluar la cadena. Repetir con un cliente válido.
\plabel{Resultado esperado} Se informa primera validación y luego la siguiente; sin ejecución de cobro.
\plabel{Resultado obtenido} CLIENTE seguido de STOCK. La integración confirma ausencia de llamada al proveedor y de persistencia parcial. El cotejo focalizado reprodujo el mismo orden con cliente ausente.
\plabel{Método y evidencia} \ruta{OperacionesComercialesTests.P23_CadenaDetienePrimeroClienteYLuegoStock}. \ruta{OperacionesComercialesTests.P23}, validación y fachada. \ruta{cotejo-p23.trx} y \ruta{P23CotejoDocumental.cs}.
\end{minipage}\par
\par\medskip\noindent\begin{minipage}{\linewidth}
\noindent\textbf{\large P24}\quad\textsf{COMPLETO}
\plabel{Descripción} Command con origen intercambiable.
\plabel{Precondiciones y entrada} Ana VIP, teletrabajo $\times$1 estándar por botón e invocador de prueba.
\plabel{Resultado esperado} Mismo comando y resultado funcional; historial de invocación.
\plabel{Resultado obtenido} ConfirmarPedidoCommand / LEGADO; 60$-$6+5=59 y auditoría de éxito.
\plabel{Método y evidencia} \ruta{OutboxComercialTests.CommandInvocaYRegistraExitoYRechazo}. \ruta{OutboxComercialTests.P24}; recorrido P24 en evidencia de fase 3 y TXT de venta.
\end{minipage}\par
\par\medskip\noindent\begin{minipage}{\linewidth}
\noindent\textbf{\large P25}\quad\textsf{COMPLETO}
\plabel{Descripción} Interpreter y coincidencia.
\plabel{Precondiciones y entrada} ESTADO = PAGADO AND TOTAL >= 50.
\plabel{Resultado esperado} Devuelve únicamente pedidos que cumplen el filtro.
\plabel{Resultado obtenido} Solo pedidos coincidentes, con pago persistido.
\plabel{Método y evidencia} \ruta{FiltroPedidosTests.P25_FiltraSobreEstadosYMontosPersistidos}. \ruta{FiltroPedidosTests.P25}; recorrido de fase 4.
\end{minipage}\par
\par\medskip\noindent\begin{minipage}{\linewidth}
\noindent\textbf{\large P26}\quad\textsf{COMPLETO}
\plabel{Descripción} Interpreter y sintaxis.
\plabel{Precondiciones y entrada} Sintaxis errónea, paréntesis y AND/OR.
\plabel{Resultado esperado} Error legible sin SQL concatenado; precedencia correcta.
\plabel{Resultado obtenido} Error con columna y precedencia; árbol tipado; consulta segura.
\plabel{Método y evidencia} \ruta{FiltroPedidosTests.P26_AndPrecedeAOrYLosParentesisCambianLaSemantica; P26_SyntaxInvalidaDaUnaColumnaYMensaje}. \ruta{FiltroPedidosTests.P26}; cierre fase 4, captura 20.
\end{minipage}\par
\par\medskip\noindent\begin{minipage}{\linewidth}
\noindent\textbf{\large P27}\quad\textsf{COMPLETO}
\plabel{Descripción} Iterator de hojas físicas.
\plabel{Precondiciones y entrada} Iterar teletrabajo y caso con ramas compartidas.
\plabel{Resultado esperado} P01, P02, P03 una vez por unidad requerida, sin bucle.
\plabel{Resultado obtenido} Tres SKU en orden determinista y multiplicadores correctos.
\plabel{Método y evidencia} \ruta{RecorridosCatalogoTests.P27_P34_ConsultaRealDelKitOficial}. \ruta{RecorridosCatalogoTests.P27_P34} e IteratorEsDeterminista.
\end{minipage}\par
\par\medskip\noindent\begin{minipage}{\linewidth}
\noindent\textbf{\large P28}\quad\textsf{COMPLETO}
\plabel{Descripción} Mediator y refresco de vistas.
\plabel{Precondiciones y entrada} Confirmar y revisar catálogo, pedidos y métricas.
\plabel{Resultado esperado} Catálogo, pedidos y métricas se refrescan sin acoplamiento directo.
\plabel{Resultado obtenido} Stock 7/4/19, PAGADO y ventas brutas recalculadas; métricas de suscriptores tras reproceso.
\plabel{Método y evidencia} \ruta{MediadorComercialTests.ConfirmacionRefrescaCincoColaboradoresSobreElEstadoConfirmado}. \ruta{MediadorComercialTests.ConfirmacionRefrescaCincoColaboradoresSobreElEstadoConfirmado}; recorrido de fase 3.
\end{minipage}\par
\par\medskip\noindent\begin{minipage}{\linewidth}
\noindent\textbf{\large P29}\quad\textsf{COMPLETO}
\plabel{Descripción} Memento y edición reversible.
\plabel{Precondiciones y entrada} Tres ediciones, dos deshacer y un rehacer, luego pagar.
\plabel{Resultado esperado} Se restauran borradores completos; no se deshace un pago.
\plabel{Resultado obtenido} Snapshots restaurados, límite cinco y pago no reversible.
\plabel{Método y evidencia} \ruta{PatronesBorradorTests.P29_DeshacerYRehacerRestauraYElPagoCierraElHistorial}. \ruta{PatronesBorradorTests.P29}; fase 4, captura 17.
\end{minipage}\par
\par\medskip\noindent\begin{minipage}{\linewidth}
\noindent\textbf{\large P30}\quad\textsf{COMPLETO}
\plabel{Descripción} Observer y deduplicación.
\plabel{Precondiciones y entrada} Tres suscriptores para compra confirmada; reintentar y provocar rollback.
\plabel{Resultado esperado} Notificación, auditoría y métrica actualizadas una vez.
\plabel{Resultado obtenido} Aviso, métrica y auditoría por suscriptor; sin duplicados ni evento exitoso de rollback.
\plabel{Método y evidencia} \ruta{OutboxComercialTests.P11_P30_CanalElegidoYTresObservadoresRecibenUnEvento; ConfirmacionTests.P20_FalloDePersistenciaReviertePagoEstadoMovimientosYEvento}. \ruta{OutboxComercialTests.P11_P30} y P37; \ruta{ConfirmacionTests.P20}.
\end{minipage}\par
\par\medskip\noindent\begin{minipage}{\linewidth}
\noindent\textbf{\large P31}\quad\textsf{COMPLETO}
\plabel{Descripción} State y transiciones.
\plabel{Precondiciones y entrada} Ciclo completo y cancelar pagado.
\plabel{Resultado esperado} Transiciones con reglas por estado, sin escritura directa.
\plabel{Resultado obtenido} Transiciones permitidas y rechazo de repetición, retroceso y cancelación pagada.
\plabel{Método y evidencia} \ruta{PedidoTests.StateRegistraSoloTransicionesAutorizadasPorElEstado}, \ruta{StatePermiteCancelarBorradorOValidado}. \ruta{PedidoTests.StateRegistraSoloTransicionesAutorizadasPorElEstado}.
\end{minipage}\par
\par\medskip\noindent\begin{minipage}{\linewidth}
\noindent\textbf{\large P32}\quad\textsf{COMPLETO}
\plabel{Descripción} Strategy y prioridad.
\plabel{Precondiciones y entrada} Promociones normal, VIP, mayorista y extensión registrada.
\plabel{Resultado esperado} Descuentos 6/15/0; extensión sin editar la calculadora central.
\plabel{Resultado obtenido} 6/15/0, un descuento, prioridad y redondeo decimal.
\plabel{Método y evidencia} \ruta{OperacionesComercialesTests.PromocionYUmbralFisico; P32_StrategyAdmiteExtensionYRedondeaMitadHaciaFuera}. \ruta{OperacionesComercialesTests.P32} y PromocionYUmbralFisico.
\end{minipage}\par
\par\medskip\noindent\begin{minipage}{\linewidth}
\noindent\textbf{\large P33}\quad\textsf{COMPLETO}
\plabel{Descripción} Template Method y plazos.
\plabel{Precondiciones y entrada} DEFECTO día 20 y ARREPENTIMIENTO día 10 en pedidos entregados.
\plabel{Resultado esperado} Defecto aceptado; arrepentimiento rechazado sin reembolso/stock.
\plabel{Resultado obtenido} Aceptación y rechazo según plazo; rechazo sin reintegro ni reembolso.
\plabel{Método y evidencia} \ruta{LogisticaTests.P33_PoliticasDePlazoYCondicion}. \ruta{LogisticaTests.P33}.
\end{minipage}\par
\par\medskip\noindent\begin{minipage}{\linewidth}
\noindent\textbf{\large P34}\quad\textsf{COMPLETO}
\plabel{Descripción} Visitor y peso logístico.
\plabel{Precondiciones y entrada} Visitar kit teletrabajo.
\plabel{Resultado esperado} Detalle P01/P02/P03; peso total 1.00 kg.
\plabel{Resultado obtenido} P01/P02/P03 y 1.00 kg; multiplicidad $\times$3 también cubierta.
\plabel{Método y evidencia} \ruta{RecorridosCatalogoTests.P27_P34_ConsultaRealDelKitOficial}. \ruta{RecorridosCatalogoTests.P27_P34}; fase 4, captura 09.
\end{minipage}\par
\subsection*{Integridad, recuperación y arquitectura}
\par\medskip\noindent\begin{minipage}{\linewidth}
\noindent\textbf{\large P35}\quad\textsf{COMPLETO}
\plabel{Descripción} Devolución y reintegro únicos.
\plabel{Precondiciones y entrada} Devolver por defecto y repetir tras reabrir.
\plabel{Resultado esperado} Para Ana: USD 54.00 de reembolso sin flete; un solo reintegro y devolución única.
\plabel{Resultado obtenido} 54.00 sin flete, una devolución, tres reintegros, rechazo de duplicado y auditoría.
\plabel{Método y evidencia} TRX focalizado y capturas de WinForms aportadas por el usuario después de confirmar y reabrir. \ruta{LogisticaTests.P35}; capturas usuario-inventario, usuario-devolucion y usuario-auditoria.
\end{minipage}\par
\par\medskip\noindent\begin{minipage}{\linewidth}
\noindent\textbf{\large P36}\quad\textsf{COMPLETO}
\plabel{Descripción} Dos cobros concurrentes.
\plabel{Precondiciones y entrada} Dos confirmaciones en conexiones distintas, stock inicial ajustado a 1.
\plabel{Resultado esperado} Máximo un cobro aprobado; stock no negativo; sin sobreventa.
\plabel{Resultado obtenido} Un pago, un evento, una llamada a pasarela y stock final 0.
\plabel{Método y evidencia} \ruta{ConfirmacionTests.P36_DosConfirmacionesCompitenPorUnaUnidad}. \ruta{ConfirmacionTests.P36}.
\end{minipage}\par
\par\medskip\noindent\begin{minipage}{\linewidth}
\noindent\textbf{\large P37}\quad\textsf{COMPLETO}
\plabel{Descripción} Reinicio y recuperación de Outbox.
\plabel{Precondiciones y entrada} Fallo de aviso, reiniciar host sobre la misma base, reintentar.
\plabel{Resultado esperado} Evento pendiente se recupera, se marca procesado una vez y sin duplicar cobro.
\plabel{Resultado obtenido} Intentos 1$\to$2, pendiente$\to$procesado, sin pago ni aviso duplicado.
\plabel{Método y evidencia} TRX focalizado, capturas originales 02/06 y consultas originales 03/08. \ruta{OutboxComercialTests.P37}; fase 4, capturas 02/06.
\end{minipage}\par
\par\medskip\noindent\begin{minipage}{\linewidth}
\noindent\textbf{\large P38}\quad\textsf{COMPLETO}
\plabel{Descripción} Build limpio y contratos.
\plabel{Precondiciones y entrada} Compilar desde artefactos nuevos y sustituir puertos y proveedor.
\plabel{Resultado esperado} Reemplazo de proveedor, contratos segregados, referencias correctas y 25 pruebas automatizadas aprobadas.
\plabel{Resultado obtenido} 119 pruebas, 0 errores/advertencias; referencias, ISP, DIP y LSP.
\plabel{Método y evidencia} TRX final, TRX focalizado y build de cierre. ContratosEscritorioTests; ReferenciasTests; \ruta{ConfirmacionTests.P16}; TRX final.
\end{minipage}\par



\subsubsection{Gestión de Git, GitHub y trazabilidad}
La interfaz, la auditoría y la documentación se consolidaron en \texttt{develop} mediante avance rápido. El historial sigue Conventional Commits y registra por separado los cambios de código, pruebas y documentación.

\begin{longtable}{@{}L{.36\linewidth}L{.17\linewidth}L{.40\linewidth}@{}}
\toprule Rama & Referencia & Alcance \\ \midrule\endhead
\texttt{main} & \texttt{9e5f254} & Base estable. \\
\texttt{develop} & \texttt{b0380bc} & Cierre académico y código consolidado. \\
\ruta{feature/interfaz-winforms} & \texttt{a4e428b} & Interfaz y corrección DPI. \\
\ruta{audit/calidad-integral} & \texttt{12458bc} & Código y evidencia de auditoría. \\
\bottomrule
\end{longtable}

\begin{longtable}{@{}L{.16\linewidth}L{.22\linewidth}L{.56\linewidth}@{}}
\toprule Commit & Tipo y ámbito & Cambio verificable \\ \midrule\endhead
\texttt{a4e428b} & fix / WinForms & Navegación y tarjetas adaptables para texto ampliado. \\
\texttt{66cd254} & fix / Mediator & Métricas refrescadas al cambiar el estado del pedido. \\
\texttt{ccc6cb3} & refactor / pedidos & Reutilizar cantidades físicas validadas. \\
\texttt{9d89b66} & test / auditoría & Aumentar evidencia de flujos comerciales, entre ellos P23, P24 y P28. \\
\texttt{12458bc} & docs / auditoría & Registrar cumplimiento, riesgos y evidencias. \\
\bottomrule
\end{longtable}

La rama \ruta{docs/revision-grafica-latex} registra la revisión gráfica desde la versión auditada. El repositorio se consulta en \url{https://github.com/Garcia2501/LogiMarket360}.

\subsubsection{Riesgos, defectos y refactorizaciones}\label{sec:riesgos}
El registro identifica 13 riesgos. R01 y R10 corresponden a refactorizaciones ejecutadas. R02--R07 y R09 registran correcciones o mitigaciones, R08 conserva una validación parcial y R11--R13 presentan propuestas sin implementar. \textbf{RNF10 sigue PARCIAL} por falta de evidencia histórica de antes y después para algunas intervenciones. Las rutas de los cuatro proyectos de aplicación parten de \ruta{src/LogiMarket.<Proyecto>}.

\par\medskip\noindent\begin{minipage}{\linewidth}\noindent\textbf{R01}\quad\textsf{Mitigado}
\plabel{Observación} Las versiones previas calculaban la valoración en \ruta{AdministrarPedidos} y \ruta{ConfirmacionPedidoFacade}. El diff elimina ambos ensamblajes.
\plabel{Consecuencia} Cotización y cobro podían divergir al cambiar tarifas o descuentos.
\plabel{Intervención} Extraer \ruta{ValorarPedido.Calcular}, \ruta{dd57d52}.
\plabel{Archivos} \ruta{Application/Precios/ValorarPedido.cs}. \ruta{Application/Pedidos/AdministrarPedidos.cs}. \ruta{Application/Pagos/ConfirmacionPedidoFacade.cs}. \ruta{Infrastructure/Persistencia/RegistroPersistencia.cs}
\plabel{Evidencia previa y posterior} Historial de fase 2 registra la selección posterior. P02–P04 y P32 volvieron a pasar. No se recuperó TRX de la ejecución anterior a la extracción.
\end{minipage}\par
\par\medskip\noindent\begin{minipage}{\linewidth}\noindent\textbf{R02}\quad\textsf{Mitigado}
\plabel{Observación} Una revocación de sesión entre autorización inicial y Proxy perdía su auditoría al revertir.
\plabel{Consecuencia} Denegación sin registro durable.
\plabel{Intervención} Drenar denegaciones pendientes fuera de la transacción revertida, \ruta{c4eb737}.
\plabel{Archivos} \ruta{Application/Seguridad/AutorizacionOperaciones.cs}. \ruta{Application/Pagos/ConfirmacionPedidoFacade.cs}. \ruta{Infrastructure/Persistencia/Repositorios/AuditoriaAcceso.cs}
\plabel{Evidencia previa y posterior} Regresión registrada en \ruta{26df60e}; \ruta{RevocacionCobroTests} vuelve a pasar.
\end{minipage}\par
\par\medskip\noindent\begin{minipage}{\linewidth}\noindent\textbf{R03}\quad\textsf{Mitigado}
\plabel{Observación} El formato compartido de moneda mostraba \ruta{USD N2} en vez de un número.
\plabel{Consecuencia} Total y confirmación ilegibles.
\plabel{Intervención} Aplicar formato \ruta{N2} al decimal y prefijar USD, \ruta{b326672}.
\plabel{Archivos} \ruta{WinForms/Presentacion.cs}
\plabel{Evidencia previa y posterior} Evidencia de fase 4 antes de la corrección; las capturas 04, 14 y las de auditoría actual muestran USD 59.00.
\end{minipage}\par
\par\medskip\noindent\begin{minipage}{\linewidth}\noindent\textbf{R04}\quad\textsf{Mitigado}
\plabel{Observación} La barra lateral reservaba una columna invisible en el acceso.
\plabel{Consecuencia} Formulario descentrado al salir o al abrir sin sesión.
\plabel{Intervención} Liberar ancho al ocultar menú, \ruta{16cf118}.
\plabel{Archivos} \ruta{WinForms/VentanaComercial.cs}
\plabel{Evidencia previa y posterior} Cierre fase 4, capturas 01 $\to$ 03.
\end{minipage}\par
\par\medskip\noindent\begin{minipage}{\linewidth}\noindent\textbf{R05}\quad\textsf{Mitigado en 96 DPI}
\plabel{Observación} A tamaño mínimo Carrito y Auditoría superponían elementos y truncaban estado e importe.
\plabel{Consecuencia} Acciones por debajo del área navegable y filas difíciles de leer.
\plabel{Intervención} Área desplazable, tamaño mínimo por columna y enfoque, \ruta{d952138}.
\plabel{Archivos} \ruta{WinForms/Presentacion.cs}. \ruta{WinForms/VentanaComercial.cs}. \ruta{WinForms/VentanaComercial.Panel.cs}
\plabel{Evidencia previa y posterior} Cierre fase 4, capturas 04–06 $\to$ 07–09 y 11.
\end{minipage}\par
\par\medskip\noindent\begin{minipage}{\linewidth}\noindent\textbf{R06}\quad\textsf{Mitigado}
\plabel{Observación} Encabezado seleccionado con contraste insuficiente.
\plabel{Consecuencia} Difícil identificar el campo activo.
\plabel{Intervención} Fondo claro para encabezados seleccionados, \ruta{bb9e06d}.
\plabel{Archivos} \ruta{WinForms/Presentacion.cs}
\plabel{Evidencia previa y posterior} Imágenes de fase 4 previas $\to$ captura 23.
\end{minipage}\par
\par\medskip\noindent\begin{minipage}{\linewidth}\noindent\textbf{R07}\quad\textsf{Mitigado}
\plabel{Observación} Acción principal deshabilitada conservaba fondo oscuro con texto inactivo.
\plabel{Consecuencia} BODEGA tenía dificultad para leer el control de cobro bloqueado.
\plabel{Intervención} Actualizar contraste en \ruta{Tema.Boton}, \ruta{3d7d51f}.
\plabel{Archivos} \ruta{WinForms/Presentacion.cs}
\plabel{Evidencia previa y posterior} Cierre 16 $\to$ 21. La prueba de autorización también confirma que el backend deniega.
\end{minipage}\par
\par\medskip\noindent\begin{minipage}{\linewidth}\noindent\textbf{R08}\quad\textsf{Parcial, bloqueante}
\plabel{Observación} Tamaños fijos de navegación y tarjetas fallaron con texto aumentado.
\plabel{Consecuencia} Rótulos y cifras recortados al 125 \% y al 150 \%.
\plabel{Intervención} Medir ancho preferido y permitir crecimiento de tarjetas, \ruta{a4e428b}.
\plabel{Archivos} \ruta{WinForms/Presentacion.cs}. \ruta{WinForms/VentanaComercial.cs}
\plabel{Evidencia previa y posterior} Capturas aportadas antes de la corrección; después se verificó 96 DPI. Faltan 120/144 DPI. La revalidación a 120 DPI amplió la evidencia; queda una transición ambigua y 144 DPI sin ejecutar tras el arreglo.
\end{minipage}\par
\par\medskip\noindent\begin{minipage}{\linewidth}\noindent\textbf{R09}\quad\textsf{Corregido}
\plabel{Observación} \ruta{FormularioPrincipalMediator.Destinos} omitía métricas en \ruta{EstadoPedido}.
\plabel{Consecuencia} Panel y contador de pendientes quedaban obsoletos tras despacho.
\plabel{Intervención} Añadir el destino de métricas a ese cambio, \ruta{66cd254}.
\plabel{Archivos} \ruta{Application/Navegacion/FormularioPrincipalMediator.cs}. \ruta{tests/Integracion/MediadorComercialTests.cs}
\plabel{Evidencia previa y posterior} \ruta{mediator-antes.trx} exige callback y recibe null $\to$ \ruta{mediator-despues.trx} 3 aprobadas.
\end{minipage}\par
\par\medskip\noindent\begin{minipage}{\linewidth}\noindent\textbf{R10}\quad\textsf{Refactorizado}
\plabel{Observación} Cadena y fachada expandían por separado el snapshot para validar y luego valorar/descontar.
\plabel{Consecuencia} Trabajo duplicado y riesgo de que las cantidades validadas diverjan de las aplicadas.
\plabel{Intervención} Compartir producto y cantidad validados, \ruta{ccc6cb3}.
\plabel{Archivos} \ruta{Application/Pedidos/CadenaValidacionPedido.cs}. \ruta{Application/Pagos/ConfirmacionPedidoFacade.cs}
\plabel{Evidencia previa y posterior} \ruta{expansion-antes.trx} y \ruta{expansion-despues.trx}, 19/19 ambos; cubre cifras, insuficiencia, rollback y concurrencia.
\end{minipage}\par
\par\medskip\noindent\begin{minipage}{\linewidth}\noindent\textbf{R11}\quad\textsf{Abierto, fuera de escala académica actual}
\plabel{Observación} \ruta{LecturaReportePedidos.ListarAsync} materializa todos los pedidos antes del árbol de filtros.
\plabel{Consecuencia} El tiempo y uso de memoria crecen con el historial.
\plabel{Intervención} Propuesta: proyección por rango y paginación, con traducción tipada que preserve precedencia. Sin commit asociado.
\plabel{Archivos} \ruta{Infrastructure/Persistencia/Repositorios/LecturaReportePedidos.cs}. \ruta{Application/Reportes/ConsultarReportePedidos.cs}
\plabel{Evidencia previa y posterior} P25/P26 prueban semántica actual. Sin medición de rendimiento a gran escala.
\end{minipage}\par
\par\medskip\noindent\begin{minipage}{\linewidth}\noindent\textbf{R12}\quad\textsf{Riesgo residual, no se reprodujo en el entorno}
\plabel{Observación} \ruta{InvocadorComandos.EjecutarAsync} espera el registro en el bloque finally. La fachada también espera trazas tras la transacción.
\plabel{Consecuencia} Un fallo de escritura diagnóstica podría ocultar el resultado comercial confirmado o la causa original.
\plabel{Intervención} Diseñar persistencia diagnóstica durable o distinguir resultado comercial de estado de auditoría. Sin commit asociado.
\plabel{Archivos} \ruta{Application/Comandos/InvocadorComandos.cs}. \ruta{Application/Pagos/ConfirmacionPedidoFacade.cs}
\plabel{Evidencia previa y posterior} Se verificaron éxito y rechazo normal en P19/P24. Sin inyección de fallo de almacenamiento después del commit.
\end{minipage}\par
\par\medskip\noindent\begin{minipage}{\linewidth}\noindent\textbf{R13}\quad\textsf{Abierto, capacidad de reproceso manual verificada}
\plabel{Observación} \ruta{ProcesadorOutboxSqlite.ProcesarAsync} selecciona por fecha un lote máximo de 100 pendientes.
\plabel{Consecuencia} Un lote bloqueado persistentemente puede retrasar eventos posteriores hasta resolver el fallo.
\plabel{Intervención} Evaluar planificación de reintentos y equidad mantenida por fecha de próximo intento. Sin commit asociado.
\plabel{Archivos} \ruta{Infrastructure/Persistencia/Eventos/ProcesadorOutboxSqlite.cs}
\plabel{Evidencia previa y posterior} P37 recupera un fallo y confirma idempotencia. No se ejecutó un backlog mayor de 100 fallos.
\end{minipage}\par




\clearpage\subsubsection{Evaluación técnica al corte}
\begin{longtable}{@{}L{.26\linewidth}L{.18\linewidth}L{.50\linewidth}@{}}
\toprule Área & Estado & Fundamento \\ \midrule\endhead
GoF & COMPLETO & 23 colaboraciones funcionales con fuente, pruebas y sus condiciones objetivas. \\
P01--P38 & COMPLETO & La matriz detallada documenta las condiciones de los 38 escenarios de la guía. \\
Integridad y seguridad & COMPLETO & Rollback, idempotencia, autorización, concurrencia, unicidad, persistencia y pruebas negativas. \\
RNF03 & PARCIAL & Cotejo visual en la sección~\ref{sec:dpi}. \\
RNF10 & PARCIAL & Evidencia por intervención en la sección~\ref{sec:riesgos}. \\
\bottomrule
\end{longtable}


\clearpage
\subsubsection{Evidencia complementaria y reproducción}
\paragraph{Comprobantes internos.}
Los archivos TXT de venta y devolución de la fase 3 complementan la evidencia de P17. Los anexos indican dónde consultar los originales y las fuentes.

\paragraph{Inventario, historial y operación de Outbox.}
La Figura~\ref{fig:inventario-venta} muestra los saldos inmediatamente posteriores a la venta, antes del reintegro que presenta la Figura~\ref{fig:inventario}. La Figura~\ref{fig:cambio-precio} registra la edición del catálogo, cuya preservación histórica se coteja con la Figura~\ref{fig:pedidos}. Ambas comparaciones se corroboran en P08 y P35.
\vistafig{inventario-venta-revision.jpg}{inventario-venta}{10.55cm}{Existencias P01/P02/P03 de 7/4/19 después del cobro y de reabrir SQLite.}
\clearpage
\vistafig{cambio-precio-revision.jpg}{cambio-precio}{10.55cm}{Precio vigente de P01 a USD 21.00 al reabrir la copia del escenario de fase 4.}

\clearpage
P21 comparó dos construcciones equivalentes de 1.000 filas con los cinco SKU de control. La Figura~\ref{fig:flyweight} recoge el recuento de objetos, identidad compartida y los bytes asignados en la ejecución. La memoria se mide para la construcción comparada, en el hilo que ejecuta la medición.
\vistafig{flyweight-cierre.jpg}{flyweight}{10.7cm}{Demostración Flyweight sobre 1.000 renglones de catálogo.}
La fábrica se crea para la proyección y usa el SKU como clave de los metadatos inmutables. Los datos extrínsecos conservan valores independientes por renglón. Este resultado caracteriza esa muestra; no estima consumo total de la aplicación ni del archivo SQLite.

\clearpage
P37 indujo un fallo del notificador y reapertura del mismo archivo. La Figura~\ref{fig:outbox-fallo} corresponde al primer intento. La Figura~\ref{fig:outbox-recuperado} muestra recuperación en el segundo intento. Los archivos \ruta{evidencias/03-sqlite-antes-de-reiniciar.json} y \ruta{evidencias/08-sqlite-recuperado.json} permiten cotejar la referencia de pago y la marca persistida de procesamiento.
\vistafig{outbox-fallo.jpg}{outbox-fallo}{7.9cm}{Evento pendiente con un intento y un error simulado registrado.}
\vistafig{outbox-recuperado.jpg}{outbox-recuperado}{7.9cm}{Evento recuperado al reprocesar y métrica de sesión registrada en Auditoría.}

\clearpage
La Figura~\ref{fig:metricas-panel} completa la comparación visual de P15 junto con la Figura~\ref{fig:outbox-recuperado}. Ambas se documentaron en el recorrido de fase 4. El contador de sesión expresa la actividad observada por el host en ejecución.
\vistafig{metricas-panel.jpg}{metricas-panel}{10.7cm}{Panel general mostrando el mismo conteo que Auditoría en el recorrido de reproceso.}

\clearpage
\paragraph{Validación negativa del filtro.}
En P26, el filtro incompleto \ruta{ESTADO =} devuelve un error con ubicación en la columna 9. La Figura~\ref{fig:interpreter-error} muestra el mensaje en la interfaz ejecutada. La prueba automatizada también confirma precedencia y manejo de paréntesis.
\vistafig{interpreter-error-cierre.jpg}{interpreter-error}{10.9cm}{Error de sintaxis localizado en el monitor de reportes y comprobantes.}


\clearpage
\paragraph{Wireframes de referencia}
Las Figuras~\ref{fig:wire-1} a \ref{fig:wire-6} reproducen los doce wireframes de las páginas 11 a 16 de la guía APE 04 \cite{ape04}. Sus indicadores y datos son esquemáticos. La implementación usa el catálogo de la Tabla~\ref{tab:catalogo-oficial} y las cifras de los casos oficiales. En particular, el rótulo P03 Cable HDMI del esquema corresponde a P03 Memoria USB en el catálogo oficial. Cada pareja se compara por sus controles y operaciones con las capturas reales del capítulo de interfaz.
\plabel{Acceso y Panel general} Controles de referencia: usuario y contraseña; indicadores y pedidos recientes.
\wirefig{wire-1.png}{wire-1}{14.5cm}{Wireframes institucionales de Acceso y Panel general. Fuente: guía APE 04, página 11.}

\clearpage
\plabel{Catálogo y Clientes} Controles de referencia: búsqueda de SKU o identificación, mantenimiento y fichas.
\wirefig{wire-2.png}{wire-2}{14.5cm}{Wireframes institucionales de Catálogo y Clientes. Fuente: guía APE 04, página 12.}

\clearpage
\plabel{Carrito y Pedidos} Controles de referencia: selección, cantidades, historial y cancelación controlada.
\wirefig{wire-3.png}{wire-3}{14.5cm}{Wireframes institucionales de Carrito y Pedidos. Fuente: guía APE 04, página 13.}

\clearpage
\plabel{Cobro y Despacho} Controles de referencia: resumen monetario, proveedor, guía y seguimiento.
\wirefig{wire-4.png}{wire-4}{14.5cm}{Wireframes institucionales de Cobro y Despacho. Fuente: guía APE 04, página 14.}

\clearpage
\plabel{Devoluciones e Inventario} Controles de referencia: política, reembolso, existencias y trazas por SKU.
\wirefig{wire-5.png}{wire-5}{14.5cm}{Wireframes institucionales de Devoluciones e Inventario. Fuente: guía APE 04, página 15.}

\clearpage
\plabel{Reportes y Auditoría} Controles de referencia: filtros, documentos, configuración simulada y eventos.
\wirefig{wire-6.png}{wire-6}{14.5cm}{Wireframes institucionales de Reportes y Auditoría. Fuente: guía APE 04, página 16.}




\clearpage\subsection{Habilidades blandas empleadas en la tarea}
\begin{itemize}
\item[$\square$] Liderazgo
\item[$\square$] Trabajo en equipo
\item[$\square$] Comunicación asertiva
\item[$\square$] La empatía
\item[$\checkmark$] Pensamiento crítico, aplicado a contrastar la guía, el código y los resultados.
\item[$\square$] Flexibilidad
\item[$\square$] La resolución de conflictos
\item[$\checkmark$] Adaptabilidad, demostrada al revalidar el diseño con cambios de tamaño y DPI.
\item[$\checkmark$] Responsabilidad, reflejada en el aislamiento de datos y el registro de verificaciones pendientes.
\item[$\checkmark$] Autogestión, mediante planificación por fases y control incremental de evidencias.
\end{itemize}

\subsection{Conclusiones}
La auditoría vinculó los 23 patrones GoF con casos de uso, código y pruebas. La suite archivada aprobó 119 pruebas. Los tres escenarios de compra y la devolución de Ana se ajustaron a las cifras oficiales. La reapertura de SQLite conservó el estado y permitió recuperar un evento sin repetir el cobro.

El expediente reúne P01--P38 y evidencia funcional de las doce vistas. El criterio transversal de presentación sigue parcial hasta comprobar físicamente los valores de Auditoría a 120 DPI y la compilación corregida a 144 DPI. La documentación se refiere al commit auditado y señala ese límite en la evaluación de entrega.

\subsection{Recomendaciones}
Repetir en Windows el salto Panel $\to$ Auditoría a 120 DPI y la matriz normal/maximizada a 144 DPI. Registrar el DPI real, el hash, los rótulos completos, los cuatro valores, el foco y la navegación. Mantener estas comprobaciones identificadas como limitaciones del alcance visual.

\clearpage\subsection{Referencias}
\begingroup
\makeatletter\renewcommand{\section}[2]{}\makeatother
\bibliographystyle{ieeetr}
\bibliography{sample}
\endgroup
\appendix\clearpage\section{Anexos}
\textbf{Cómo localizar los componentes del expediente.}\par
\begin{description}
\item[Fuentes] \ruta{main.tex} contiene la estructura institucional. Los bloques técnicos se incluyen desde \ruta{resultados/}. La bibliografía utiliza \ruta{sample.bib}.
\item[Diagramas y código] \ruta{diagramas/} conserva las fuentes TikZ editables. \ruta{codigo/} contiene los fragmentos citados con numeración respecto al commit auditado.
\item[Evidencia visual] \ruta{capturas/} conserva las capturas de ejecución y su manifiesto en \ruta{evidencias/capturas.md}. Los wireframes, fuentes de referencia de diseño, se guardan en \ruta{wireframes/}.
\item[Evidencia automatizada] \ruta{evidencias/auditoria-final.trx} contiene la ejecución de auditoría; se incluyen también las comparaciones de regresión y el registro de SQLite nueva.
\item[Estado de entrega] Se conservaron siete capturas auténticas con cursor o halo, sin alterar datos ni resultados funcionales. El registro de cierre visual documenta su procedencia y el alcance de las comprobaciones DPI.
\end{description}

\end{document}


